% ============================================================
%  LiQuID User Guide Manual
%  Lindblad Quantum Integrated Dynamics — v0.7.0
%  Production-grade enterprise documentation
% ============================================================
\documentclass[11pt,a4paper,twoside]{book}

% -- Encoding & fonts -----------------------------------------
\usepackage[T1]{fontenc}
\usepackage[utf8]{inputenc}
\usepackage{lmodern}
\usepackage{microtype}

% -- Page geometry --------------------------------------------
\usepackage[
  a4paper,
  top=25mm, bottom=30mm,
  inner=30mm, outer=25mm,
  headsep=8mm, footskip=12mm,
  marginparwidth=18mm
]{geometry}

% -- Colours --------------------------------------------------
\usepackage[dvipsnames,svgnames,x11names]{xcolor}
\definecolor{LiquidBlue}   {HTML}{1A3A5C}
\definecolor{LiquidTeal}   {HTML}{0A7B87}
\definecolor{LiquidGold}   {HTML}{C8922A}
\definecolor{LiquidRed}    {HTML}{B93A2B}
\definecolor{LiquidGreen}  {HTML}{276B3D}
\definecolor{LiquidGray}   {HTML}{4A4A4A}
\definecolor{LiquidLtGray} {HTML}{F4F5F7}
\definecolor{LiquidMid}    {HTML}{D0D5DE}
\definecolor{CodeBg}       {HTML}{F7F8FA}
\definecolor{CodeFrame}    {HTML}{CBD0DA}
\definecolor{NoteBlue}     {HTML}{E8F3FB}
\definecolor{NoteBlueBorder}{HTML}{2F7EBD}
\definecolor{WarnYellow}   {HTML}{FFF8E1}
\definecolor{WarnBorder}   {HTML}{D4860A}
\definecolor{PerfGreen}    {HTML}{EBF8EF}
\definecolor{PerfBorder}   {HTML}{2E8B57}
\definecolor{DangerRed}    {HTML}{FDECEA}
\definecolor{DangerBorder} {HTML}{B93A2B}

% -- Mathematics ----------------------------------------------
\usepackage{amsmath,amssymb,amsthm,mathtools,bm}
\usepackage{physics}   % \ket, \bra, \braket, \norm etc.

% -- Source code (minted) -------------------------------------
\usepackage[
  chapter,
  newfloat
]{minted}
\setminted{
  style        = tango,
  bgcolor      = CodeBg,
  fontsize     = \small,
  breaklines   = true,
  breakanywhere= false,
  autogobble   = true,
  tabsize      = 4,
  frame        = leftline,
  framesep     = 6pt,
  rulecolor    = CodeFrame,
  numbersep    = 8pt,
  xleftmargin  = 12pt,
}
\setminted[cpp]{
  style        = tango,
  bgcolor      = CodeBg,
}
\setminted[bash]{
  style        = tango,
  bgcolor      = CodeBg,
  fontsize     = \small,
}
\setminted[cmake]{
  style        = friendly,
  bgcolor      = CodeBg,
}
\setminted[python]{
  style        = friendly,
  bgcolor      = CodeBg,
}

% -- Callout boxes (tcolorbox) ---------------------------------
\usepackage[skins,breakable,listings]{tcolorbox}

%  NOTE
\newtcolorbox{notebox}[1][]{
  enhanced, breakable,
  colback   = NoteBlue,
  colframe  = NoteBlueBorder,
  coltext   = LiquidGray,
  arc       = 2pt,
  boxrule   = 1.5pt,
  left      = 8pt, right = 8pt, top = 6pt, bottom = 6pt,
  fonttitle = \bfseries\small\color{NoteBlueBorder},
  title     = {\faInfoCircle\quad NOTE},
  #1
}

%  WARNING
\newtcolorbox{warnbox}[1][]{
  enhanced, breakable,
  colback   = WarnYellow,
  colframe  = WarnBorder,
  coltext   = LiquidGray,
  arc       = 2pt,
  boxrule   = 1.5pt,
  left      = 8pt, right = 8pt, top = 6pt, bottom = 6pt,
  fonttitle = \bfseries\small\color{WarnBorder},
  title     = {\faExclamationTriangle\quad WARNING},
  #1
}

%  PERFORMANCE TIP
\newtcolorbox{perfbox}[1][]{
  enhanced, breakable,
  colback   = PerfGreen,
  colframe  = PerfBorder,
  coltext   = LiquidGray,
  arc       = 2pt,
  boxrule   = 1.5pt,
  left      = 8pt, right = 8pt, top = 6pt, bottom = 6pt,
  fonttitle = \bfseries\small\color{PerfBorder},
  title     = {\faBolt\quad PERFORMANCE TIP},
  #1
}

%  INVARIANT
\newtcolorbox{invbox}[1][]{
  enhanced, breakable,
  colback   = DangerRed,
  colframe  = DangerBorder,
  coltext   = LiquidGray,
  arc       = 2pt,
  boxrule   = 1.5pt,
  left      = 8pt, right = 8pt, top = 6pt, bottom = 6pt,
  fonttitle = \bfseries\small\color{DangerBorder},
  title     = {\faLock\quad INVARIANT},
  #1
}

% Icons (text fallbacks since fontawesome5 may not be installed)
% We redefine these as text macros for portability
\newcommand{\faInfoCircle}{[\textbf{i}]}
\newcommand{\faExclamationTriangle}{[!]}
\newcommand{\faBolt}{[\textbf{P}]}
\newcommand{\faLock}{[\textbf{INV}]}

% -- Section styling -------------------------------------------
\usepackage{titlesec}
\titleformat{\chapter}[display]
  {\normalfont\huge\bfseries\color{LiquidBlue}}
  {\normalfont\Large\color{LiquidTeal}\chaptername\ \thechapter}
  {10pt}
  {\Huge\bfseries\color{LiquidBlue}}
  [\vspace{2pt}{\color{LiquidGold}\rule{\linewidth}{2pt}}]

\titleformat{\section}
  {\normalfont\Large\bfseries\color{LiquidBlue}}
  {\thesection}{1em}{}
  [\vspace{1pt}{\color{LiquidMid}\rule{\linewidth}{0.4pt}}]

\titleformat{\subsection}
  {\normalfont\large\bfseries\color{LiquidGray}}
  {\thesubsection}{1em}{}

\titleformat{\subsubsection}
  {\normalfont\normalsize\bfseries\itshape\color{LiquidGray}}
  {\thesubsubsection}{1em}{}

% -- Header / Footer -------------------------------------------
\usepackage{fancyhdr}
\pagestyle{fancy}
\fancyhf{}
\fancyhead[LE]{\small\color{LiquidGray}\textit{\leftmark}}
\fancyhead[RO]{\small\color{LiquidGray}\textit{\rightmark}}
\fancyhead[RE,LO]{\small\color{LiquidTeal}\textbf{LiQuID v0.7.0}}
\fancyfoot[C]{\small\color{LiquidGray}\thepage}
\renewcommand{\headrulewidth}{0.4pt}
\renewcommand{\headrule}{\hbox to\headwidth{\color{LiquidMid}\leaders\hrule\hfil}}

% -- Tables ---------------------------------------------------
\usepackage{booktabs,array,tabularx,longtable}
\newcolumntype{L}[1]{>{\raggedright\arraybackslash}p{#1}}
\newcolumntype{C}[1]{>{\centering\arraybackslash}p{#1}}

% -- Lists ----------------------------------------------------
\usepackage{enumitem}
\setlist[itemize]{leftmargin=1.5em,itemsep=2pt}
\setlist[enumerate]{leftmargin=1.5em,itemsep=2pt}

% -- Figures --------------------------------------------------
\usepackage{graphicx,caption,subcaption}
\usepackage{tikz,pgfplots}
\pgfplotsset{compat=1.18}
\usetikzlibrary{arrows.meta,positioning,shapes.geometric,fit,backgrounds,
                calc,decorations.pathmorphing,matrix,chains}

% -- Cross-references & hyperlinks -----------------------------
\usepackage[
  colorlinks = true,
  linkcolor  = LiquidBlue,
  citecolor  = LiquidTeal,
  urlcolor   = LiquidTeal,
  pdftitle   = {LiQuID User Guide v0.7.0},
  pdfauthor  = {LiQuID Project},
  pdfsubject = {Monte Carlo Wave Function simulation in C++17},
  bookmarksnumbered = true,
  bookmarksopen     = true
]{hyperref}
\usepackage{cleveref}

% -- Bibliography ---------------------------------------------
\usepackage[numbers,sort&compress]{natbib}

% -- Miscellaneous ---------------------------------------------
\usepackage{parskip}
\usepackage{mdframed}
\usepackage{multicol}
\usepackage{doi}
\usepackage{url}
\usepackage{float}

% -- Custom math operators -------------------------------------
\DeclareMathOperator{\lqTr}{Tr}
\DeclareMathOperator{\lqRe}{Re}
\DeclareMathOperator{\lqIm}{Im}
\newcommand{\Heff}{H_{\mathrm{eff}}}
\newcommand{\HeffDef}{H - \frac{i}{2}\sum_k L_k^\dagger L_k}
\newcommand{\densmat}{\hat{\varrho}}    % density matrix
\newcommand{\Lindblad}{\mathcal{L}}  % Lindblad superoperator
\newcommand{\Hilbert}{\mathcal{H}}   % Hilbert space
\newcommand{\Order}{\mathcal{O}}
\newcommand{\complex}{\mathbb{C}}
\newcommand{\reals}{\mathbb{R}}
\newcommand{\liquid}{\texttt{LiQuID}}
\newcommand{\code}[1]{\texttt{\small#1}}
\newcommand{\filename}[1]{\texttt{#1}}

% -- Theorem environments --------------------------------------
\theoremstyle{definition}
\newtheorem{definition}{Definition}[chapter]
\newtheorem{algorithm}{Algorithm}[chapter]
\theoremstyle{remark}
\newtheorem{remark}{Remark}[chapter]

% -- Suppress paragraph indent (parskip handles spacing) -------
\setlength{\parindent}{0pt}

% -------------------------------------------------------------
%  DOCUMENT
% -------------------------------------------------------------
\begin{document}

% -- Title page -----------------------------------------------
\begin{titlepage}
\pagecolor{LiquidBlue}
\color{white}
\vspace*{20mm}
\begin{center}
  {\fontsize{52}{60}\selectfont\bfseries LiQuID}\\[6pt]
  {\Large\color{LiquidGold} Lindblad Quantum Integrated Dynamics}\\[4mm]
  {\large Version 0.7.0}\\[20mm]
  {\color{LiquidMid}\rule{0.7\linewidth}{1pt}}\\[8mm]
  {\LARGE\bfseries User Guide Manual}\\[6mm]
  {\large A Comprehensive Reference for Physicists and Quantum Software Engineers}\\[12mm]
  {\color{LiquidMid}\rule{0.7\linewidth}{1pt}}\\[16mm]
  {\normalsize
    Monte Carlo Wave Function Simulation $\cdot$
    Adaptive ODE Solvers $\cdot$
    Sparse Operators $\cdot$
    Parallel Ensemble Management
  }\\[8mm]
  {\normalsize C++17 $\cdot$ OpenMP $\cdot$ Header-Only Core}\\[20mm]
  {\small\color{LiquidMid} Compiled \today}
\end{center}
\vfill
\begin{center}
  {\small\color{LiquidMid}
    This document covers the complete public API, mathematical foundations,
    performance model, and extension seams of LiQuID v0.7.0.
  }
\end{center}
\end{titlepage}
\nopagecolor
\color{black}

% -- Front matter ---------------------------------------------
\frontmatter
\pagestyle{fancy}

% -------------------------------------------------------------
%  PREFACE
% -------------------------------------------------------------
\chapter*{Preface}
\addcontentsline{toc}{chapter}{Preface}

\liquid{} (\emph{Lindblad Quantum Integrated Dynamics}) is a C++17 library for
high-performance simulation of open quantum systems via the
\emph{Monte Carlo Wave Function} (MCWF) method. It is designed for research
physicists, quantum optics experimentalists, and quantum software engineers who
require both mathematical rigour and production-grade performance in a single,
dependency-free framework.

\medskip
\textbf{What this manual covers.}
This guide documents every public class, function, and configuration parameter
in LiQuID v0.7.0. It is organized in layers that mirror the library's own
dependency graph: from foundational type aliases and the random number generator,
through the linear algebra primitives and ODE solvers, to the complete
simulation pipeline. Each section anchors code examples directly to the
underlying mathematics and explains the memory management decisions that
determine performance on modern hardware.

\medskip
\textbf{What this manual does not cover.}
Density-matrix (Lindblad master equation) solvers, GPU backends, and
Python bindings are explicitly out of scope for v0.7.0. MCWF is the
single computational paradigm of the library; this focus is intentional
and is discussed in \cref{chap:theory}.

\medskip
\textbf{How to use this manual.}
Readers new to the MCWF method should start with \cref{chap:theory} before
examining the API. Experienced users may proceed directly to \cref{chap:quickstart}.
The performance chapter (\cref{chap:performance}) should be consulted before
deploying simulations with $N \geq 100$ or more than $10^4$ trajectories.

% -------------------------------------------------------------
%  STYLE GUIDE COMPONENT LIBRARY
% -------------------------------------------------------------
\chapter*{Documentation Style Guide}
\addcontentsline{toc}{chapter}{Documentation Style Guide}

This chapter defines the visual language used throughout the manual.
All code examples, mathematical formulations, and callout blocks follow
the conventions below without exception.

\section*{C++ Code Blocks}

All C++ listings use \texttt{minted} with the \texttt{tango} style and the
\texttt{cpp} lexer. Line numbers are shown for blocks exceeding ten lines.
Identifiers are cross-referenced to their defining headers on first use.

Example of a fully annotated, non-abbreviated code block:

\begin{minted}[linenos]{cpp}
#include "liquid/liquid.hpp"
using namespace liquid;

// Build the system: two-level atom, decay rate gamma
DenseOperator make_sigma_minus(double gamma) {
    DenseOperator L(2);               // 2x2 zero-initialized
    L(1, 0) = Scalar{std::sqrt(gamma), 0.0};  // sigma_- = |0><1|
    return L;
}
\end{minted}

Inline code references are typeset as \code{liquid::StateVector} or
\code{apply\_add()}.

\section*{Mathematical Formulations}

All equations are typeset in display mode with \texttt{amsmath}.
Quantum states use Dirac notation via the \texttt{physics} package.
The Lindblad master equation is the central object and is always written as:

\begin{equation}
  \frac{d\densmat}{dt}
    = -i\bigl[H, \densmat\bigr]
      + \sum_k \left(
          L_k \densmat L_k^\dagger
          - \frac{1}{2}\bigl\{L_k^\dagger L_k,\, \densmat\bigr\}
        \right)
  \label{eq:lindblad-master}
\end{equation}

The effective non-Hermitian Hamiltonian used by the MCWF algorithm is:
\begin{equation}
  \Heff = \HeffDef
  \label{eq:heff-def}
\end{equation}

\section*{Callout Block Reference}

\begin{notebox}
  A \textbf{NOTE} provides supplementary context, historical background,
  or a clarification of non-obvious behaviour.
\end{notebox}

\begin{warnbox}
  A \textbf{WARNING} identifies a precondition that, if violated, produces
  incorrect results, undefined behaviour, or numerical failure. Always read
  warnings before using the associated API.
\end{warnbox}

\begin{perfbox}
  A \textbf{PERFORMANCE TIP} describes a measurable optimisation strategy
  derived from the memory layout or computational complexity analysis
  presented in \cref{chap:performance}.
\end{perfbox}

\begin{invbox}
  An \textbf{INVARIANT} is a mathematical or memory constraint that
  LiQuID guarantees at all times. Violating an invariant from outside
  the library constitutes undefined behaviour.
\end{invbox}

% -------------------------------------------------------------
%  TABLE OF CONTENTS
% -------------------------------------------------------------
\tableofcontents
\listoffigures
\listoftables

% -------------------------------------------------------------
%  MAIN MATTER
% -------------------------------------------------------------
\mainmatter

% ============================================================
%  PART I — FOUNDATIONS
% ============================================================
\part{Foundations}

% -------------------------------------------------------------
\chapter{Introduction and Architecture}
\label{chap:intro}
% -------------------------------------------------------------

\section{Scope and Identity}
\label{sec:scope}

\liquid{} is not a general-purpose quantum computing toolkit.
Its scope is precisely stated:

\begin{quotation}\itshape
  LiQuID provides production-grade Monte Carlo Wave Function simulation
  of Markovian open quantum systems described by the Lindblad master
  equation, with adaptive trajectory management, convergence-based
  stopping, and an extensible ensemble allocation policy.
\end{quotation}

The library deliberately excludes density-matrix evolution, non-Markovian
dynamics, and circuit-model quantum computation. This focus ensures that
every design decision — from the choice of \code{Scalar = std::complex<double>}
to the FSAL property of the Dormand--Prince integrator — is motivated by
a single, well-understood computational task.

\section{Repository Structure}
\label{sec:repo-structure}

\liquid{} is distributed as a single static library with a header-only core.
The \emph{only} compiled translation unit is \filename{src/core/rng.cpp},
which implements the xoshiro256** random number generator.

\begin{figure}[H]
\centering
\begin{tikzpicture}[
  font=\small,
  box/.style={rectangle, rounded corners=3pt, draw=LiquidMid,
              fill=LiquidLtGray, text width=38mm, align=center,
              minimum height=8mm, inner sep=4pt},
  layer/.style={rectangle, rounded corners=3pt, draw=LiquidBlue,
                fill=LiquidBlue!10, text width=120mm, align=left,
                minimum height=12mm, inner sep=6pt},
  arr/.style={-{Stealth[length=4pt]}, thick, color=LiquidTeal},
]
% Layers
\node[layer] (l5) at (0, 0)
  {\textbf{Layer 5 — Simulation Interface}
   \quad\small\texttt{simulation.hpp}\quad\texttt{parameter\_sweep.hpp}};
\node[layer, below=5mm of l5] (l4)
  {\textbf{Layer 4 — Ensemble Management}
   \quad\small\texttt{ensemble\_manager.hpp}\quad\texttt{adaptive\_allocator.hpp}};
\node[layer, below=5mm of l4] (l3)
  {\textbf{Layer 3 — Trajectory}
   \quad\small\texttt{mcwf\_propagator.hpp}\quad\texttt{trajectory\_state.hpp}};
\node[layer, below=5mm of l3] (l2)
  {\textbf{Layer 2 — Quantum System}
   \quad\small\texttt{open\_system.hpp}\quad\texttt{lindblad.hpp}};
\node[layer, below=5mm of l2] (l1)
  {\textbf{Layer 1 — Linear Algebra / ODE}
   \quad\small\texttt{dense.hpp}\quad\texttt{sparse.hpp}\quad\texttt{rk4.hpp}\quad\texttt{dopri45.hpp}};
\node[layer, below=5mm of l1] (l0)
  {\textbf{Layer 0 — Foundation}
   \quad\small\texttt{types.hpp}\quad\texttt{config.hpp}\quad\texttt{rng.hpp}};

\draw[arr] (l5) -- (l4);
\draw[arr] (l4) -- (l3);
\draw[arr] (l3) -- (l2);
\draw[arr] (l2) -- (l1);
\draw[arr] (l1) -- (l0);
\end{tikzpicture}
\caption{LiQuID dependency graph. Arrows point from dependent to dependency.
         All dependencies are strictly unidirectional; no circular includes exist.}
\label{fig:layer-graph}
\end{figure}

\section{Build System and Requirements}
\label{sec:build}

\subsection{Prerequisites}
\label{subsec:prerequisites}

\begin{table}[H]
\centering
\caption{Build requirements for LiQuID v0.7.0.}
\label{tab:requirements}
\begin{tabularx}{\linewidth}{L{30mm}L{25mm}X}
\toprule
\textbf{Component} & \textbf{Version} & \textbf{Notes} \\
\midrule
CMake      & $\geq$ 3.16   & Required for \code{target\_compile\_features} \\
C++ compiler & C++17       & GCC $\geq$ 7, Clang $\geq$ 5, MSVC 2019+ \\
OpenMP     & Optional       & Enables multi-threaded ensemble dispatch \\
Python + Pygments & Optional & Required only for \texttt{minted}-rendered docs \\
\bottomrule
\end{tabularx}
\end{table}

\subsection{CMake Configuration}
\label{subsec:cmake}

\begin{minted}[linenos]{bash}
git clone https://github.com/your-org/LiQuID.git
cd LiQuID
cmake -B build -DCMAKE_BUILD_TYPE=Release \
               -DLIQUID_OPENMP=ON          \
               -DLIQUID_TESTS=ON           \
               -DLIQUID_EXAMPLES=ON
cmake --build build --parallel
ctest --test-dir build --output-on-failure
\end{minted}

CMake exposes four boolean options:

\begin{table}[H]
\centering
\caption{CMake build options.}
\label{tab:cmake-options}
\begin{tabularx}{\linewidth}{L{42mm}L{18mm}X}
\toprule
\textbf{Option} & \textbf{Default} & \textbf{Effect} \\
\midrule
\code{LIQUID\_OPENMP}    & \code{ON}  & Links OpenMP if found; falls back silently \\
\code{LIQUID\_TESTS}     & \code{ON}  & Builds 10 unit + 7 validation test executables \\
\code{LIQUID\_EXAMPLES}  & \code{ON}  & Builds \code{jc\_steady\_state} and \code{single} \\
\code{LIQUID\_SANITIZE}  & \code{OFF} & Enables AddressSanitizer + UBSan (Debug only) \\
\bottomrule
\end{tabularx}
\end{table}

\subsection{Linking Against LiQuID}
\label{subsec:linking}

\begin{minted}{cmake}
find_package(LiQuID REQUIRED)
target_link_libraries(my_simulation PRIVATE liquid)
\end{minted}

Alternatively, add as a subdirectory:

\begin{minted}{cmake}
add_subdirectory(extern/LiQuID)
target_link_libraries(my_simulation PRIVATE liquid)
\end{minted}

\begin{notebox}
  The umbrella header \code{liquid/liquid.hpp} includes all public
  headers in correct dependency order. For large projects, you may
  include individual layer headers to reduce compilation time.
\end{notebox}

% -------------------------------------------------------------
\chapter{Theoretical Foundations}
\label{chap:theory}
% -------------------------------------------------------------

\section{The Lindblad Master Equation}
\label{sec:lindblad-master}

An open quantum system is one that interacts with an unmonitored
environment (reservoir). When the system--reservoir coupling is weak
(\emph{Born approximation}) and the reservoir correlation time is much
shorter than the system dynamics (\emph{Markov approximation}), the
reduced density matrix $\densmat$ of the system evolves under the
\emph{Gorini--Kossakowski--Sudarshan--Lindblad} (GKSL) master equation:

\begin{equation}
  \boxed{
  \frac{d\densmat}{dt}
    = -i\bigl[H, \densmat\bigr]
      + \sum_{k=1}^{K} \left(
          L_k \densmat L_k^\dagger
          - \frac{1}{2} L_k^\dagger L_k \densmat
          - \frac{1}{2} \densmat\, L_k^\dagger L_k
        \right)
  }
  \label{eq:lindblad}
\end{equation}

where:
\begin{itemize}
  \item $H \in \complex^{N \times N}$ is the Hermitian system Hamiltonian
        (\emph{physical units}: $\hbar = 1$);
  \item $\{L_k\}_{k=1}^{K}$ are the \emph{Lindblad jump operators}, which
        encode the dissipative channels of the system--environment interaction;
  \item $[\cdot,\cdot]$ denotes the commutator.
\end{itemize}

\cref{eq:lindblad} is the most general Markovian master equation that
preserves the trace, Hermiticity, and positive semi-definiteness of
$\densmat$.

\begin{notebox}
  \liquid{} works in natural units $\hbar = 1$. All energies and rates
  are expressed in the same units as the Hamiltonian matrix elements.
  Concretely, if $H_{00}$ has units of MHz, then decay rates $\gamma_k$
  must also be in MHz, and $t_{\mathrm{final}}$ in $\mu\mathrm{s}^{-1}$.
\end{notebox}

\subsection{Physical Interpretation of Jump Operators}
\label{subsec:jump-ops}

Each operator $L_k$ describes one incoherent process. Common examples
in quantum optics are:

\begin{table}[H]
\centering
\caption{Standard Lindblad jump operators for common dissipation channels.}
\label{tab:jump-ops}
\begin{tabularx}{\linewidth}{L{38mm}L{38mm}X}
\toprule
\textbf{Process} & \textbf{Operator} & \textbf{C++ construction} \\
\midrule
Spontaneous emission (rate $\gamma$)
  & $L = \sqrt{\gamma}\,\sigma_-$
  & \code{L(1,0) = Scalar\{sqrt(gamma), 0\};} \\[4pt]
Cavity photon loss (rate $\kappa$)
  & $L = \sqrt{\kappa}\,a$
  & Lower-triangular with $\sqrt{n}$ entries \\[4pt]
Pure dephasing (rate $\gamma_\phi$)
  & $L = \sqrt{\gamma_\phi}\,\sigma_z$
  & Diagonal: $\pm\sqrt{\gamma_\phi}$ \\[4pt]
Thermal excitation (rate $\gamma\bar{n}$)
  & $L = \sqrt{\gamma\bar{n}}\,\sigma_+$
  & \code{L(0,1) = Scalar\{sqrt(gamma\_n), 0\};} \\
\bottomrule
\end{tabularx}
\end{table}

\section{The MCWF Stochastic Unravelling}
\label{sec:mcwf-theory}

\subsection{Core Idea}
\label{subsec:mcwf-core}

The MCWF method (Dalibard, Castin and M{\o}lmer, 1992; Dum, Zoller and Ritsch,
1992) \emph{unravels} the master equation into an ensemble of pure-state
trajectories. Each trajectory is a stochastic process on the Hilbert space
$\Hilbert = \complex^N$, and ensemble averages reproduce the density matrix:

\begin{equation}
  \densmat(t) = \overline{\ket{\psi(t)}\bra{\psi(t)}}
  \label{eq:mcwf-avg}
\end{equation}

where the overline denotes an average over many independent realizations
of the stochastic process. The computational advantage over direct
master-equation integration is memory: storing $\ket{\psi} \in \complex^N$
requires $\Order(N)$ complex numbers, whereas $\densmat$ requires
$\Order(N^2)$. For a 100-photon Fock space ($N = 101$), this is the
difference between $\sim 1.6\,\mathrm{kB}$ and $\sim 160\,\mathrm{kB}$.
For $N = 1000$, the ratio grows to $\times 1000$.

\subsection{The Effective Non-Hermitian Hamiltonian}
\label{subsec:heff}

Define the aggregate decay operator:
\begin{equation}
  \Gamma \equiv \sum_{k=1}^{K} L_k^\dagger L_k
  \label{eq:Gamma-def}
\end{equation}

and the effective Hamiltonian:
\begin{equation}
  \Heff \equiv H - \frac{i}{2}\,\Gamma = \HeffDef
  \label{eq:heff}
\end{equation}

$\Heff$ is \emph{not Hermitian}: $\Heff \neq \Heff^\dagger$.
Its non-Hermitian part $-\tfrac{i}{2}\Gamma$ causes the norm of
$\ket{\psi}$ to decrease over time, encoding the probability of
\emph{no jump having occurred} up to time $t$.

\begin{invbox}
  \textbf{INV-$\Gamma$:} In LiQuID, $\Gamma = \sum_k L_k^\dagger L_k$ is
  precomputed and stored immutably in \code{LindbladSet::gamma\_} at
  construction time. It is never recomputed. Modifying the jump operators
  after construction is not possible because \code{LindbladSet::ops\_}
  is private and the copy constructor is deleted.
\end{invbox}

\subsection{Between-Jump Evolution}
\label{subsec:between-jump}

Between quantum jumps, the wavefunction satisfies the
\emph{non-Hermitian Schr{\"o}dinger equation}:

\begin{equation}
  \frac{d\ket{\psi}}{dt} = -i\,\Heff\ket{\psi}
  \label{eq:nhse}
\end{equation}

This is integrated numerically by the ODE solver.
The norm $\|\ket{\psi(t)}\|^2 = \braket{\psi(t)|\psi(t)}$ decreases
monotonically between jumps. Specifically, at time $t$ after the last
normalization event:

\begin{equation}
  \frac{d}{dt}\|\ket{\psi}\|^2
    = -\braket{\psi|\Gamma|\psi}
    = -\sum_k \|L_k\ket{\psi}\|^2 \leq 0
  \label{eq:norm-decay}
\end{equation}

\subsection{Jump Detection and Channel Selection}
\label{subsec:jump-detection}

A random threshold $r \sim \mathrm{Uniform}(0, 1]$ is drawn at the
start of each inter-jump interval (and redrawn immediately after every
jump). A jump is detected when:

\begin{equation}
  \|\ket{\psi(t)}\|^2 < r
  \label{eq:jump-detect}
\end{equation}

Upon detection, the jump channel $k^*$ is selected with probability
proportional to $\|L_k\ket{\psi}\|^2$:

\begin{equation}
  \Pr(\text{channel } k) = \frac{\|L_k\ket{\psi}\|^2}{\sum_{j}\|L_j\ket{\psi}\|^2}
  \label{eq:jump-prob}
\end{equation}

The post-jump state is renormalized:

\begin{equation}
  \ket{\psi} \leftarrow \frac{L_{k^*}\ket{\psi}}{\|L_{k^*}\ket{\psi}\|}
  \label{eq:post-jump}
\end{equation}

\subsection{Observable Estimation}
\label{subsec:observables}

For a Hermitian observable $O$, the MCWF estimator is:

\begin{equation}
  \langle O \rangle(t)
    \approx \frac{1}{M} \sum_{m=1}^{M}
      \frac{\braket{\psi^{(m)}(t)|O|\psi^{(m)}(t)}}
           {\braket{\psi^{(m)}(t)|\psi^{(m)}(t)}}
  \label{eq:observable-est}
\end{equation}

where the denominator accounts for the fact that trajectories are stored
in an unnormalized form between output intervals. In LiQuID, the
\code{ObservableFn} callable is responsible for the normalization;
the \code{ObservableAccumulator} computes only the sum and variance.

\subsection{Algorithm Summary}
\label{subsec:algorithm-summary}

\begin{algorithm}[H]
\label{alg:mcwf}
\textbf{Algorithm: MCWF Single Trajectory} (one call to
\code{MCWFPropagator::advance})

\begin{enumerate}[leftmargin=2em,label=\arabic*.]
  \item Select step size $\delta t$ (warm-start from previous step).
  \item \textbf{ODE step}: integrate \cref{eq:nhse} from $t$ to $t + \delta t$
        using the configured stepper (RK4 or DOPRI45), yielding
        $\ket{\psi_{\mathrm{prop}}}$.
  \item \textbf{Jump check}: if $\|\ket{\psi_{\mathrm{prop}}}\|^2 > r$:
    \begin{enumerate}[label=3\alph*.]
      \item Accept step: $\ket{\psi} \leftarrow \ket{\psi_{\mathrm{prop}}}$,
            $t \leftarrow t + \delta t$.
      \item Return \code{Advanced} (or \code{Completed} if $t \geq t_{\mathrm{final}}$).
    \end{enumerate}
  \item \textbf{Jump handling}: a jump occurred in $[t, t+\delta t]$:
    \begin{enumerate}[label=4\alph*.]
      \item (Phase 3) Locate exact jump time via bisection.
      \item Compute $p_k = \|L_k\ket{\psi}\|^2$ for all $k$.
      \item Draw channel $k^* \sim \{p_k / \sum p_k\}$.
      \item Apply: $\ket{\psi} \leftarrow L_{k^*}\ket{\psi} / \|L_{k^*}\ket{\psi}\|$.
      \item Redraw $r \sim \mathrm{Uniform}(0,1]$.
      \item Reset stepper (state is discontinuous).
      \item Return \code{JumpOccurred}.
    \end{enumerate}
  \item If numerical failure (step size collapse): return \code{Failed}.
\end{enumerate}
\end{algorithm}

\begin{notebox}
  \code{advance()} returns at the \emph{first} jump, not after all jumps
  to $t_{\mathrm{target}}$. This event-driven design allows the
  \code{EnsembleManager} to inspect, suspend, or redirect a trajectory
  after each jump without re-entering the propagation loop.
\end{notebox}

% ============================================================
%  PART II — API REFERENCE
% ============================================================
\part{API Reference}

% -------------------------------------------------------------
\chapter{Quick Start}
\label{chap:quickstart}
% -------------------------------------------------------------

\section{The Five-Minute Example}
\label{sec:five-minute}

The following complete program simulates a two-level atom decaying
from its excited state at rate $\gamma = 1.0$, with a resonance
frequency $\omega = 1.0$. The exact solution for the expectation value
of $\sigma_z = \ket{e}\bra{e} - \ket{g}\bra{g}$ is:

\begin{equation}
  \langle \sigma_z(t) \rangle = 2 e^{-\gamma t} - 1
  \label{eq:two-level-exact}
\end{equation}

\begin{minted}[linenos]{cpp}
#include "liquid/liquid.hpp"
using namespace liquid;
using namespace liquid::ensemble;

int main() {
    // -- Physical parameters ------------------------------------------------
    const double omega = 1.0;    // transition frequency [arbitrary units]
    const double gamma = 1.0;    // spontaneous emission rate
    const double T     = 3.0;    // total simulation time

    // -- Hamiltonian: H = (omega/2) * sigma_z ------------------------------
    // Basis: |0> = excited, |1> = ground
    DenseOperator H(2);
    H(0, 0) = Scalar{ omega / 2.0, 0.0};
    H(1, 1) = Scalar{-omega / 2.0, 0.0};

    // -- Jump operator: L = sqrt(gamma) * sigma_minus ----------------------
    // sigma_minus = |1><0|: maps excited |0> to ground |1>
    DenseOperator L(2);
    L(1, 0) = Scalar{std::sqrt(gamma), 0.0};

    // -- Observable: sigma_z = |0><0| - |1><1| -----------------------------
    DenseOperator sz(2);
    sz(0, 0) = Scalar{ 1.0, 0.0};
    sz(1, 1) = Scalar{-1.0, 0.0};

    // -- Build simulation --------------------------------------------------
    Simulation sim = SimulationBuilder{}
        .hamiltonian(std::move(H))
        .collapse_operator(std::move(L))
        .observe("sigma_z", std::move(sz))
        .seed(42)
        .dt(1e-3)
        .stop_at_sem(0.01)          // stop when relative SEM < 1%
        .min_trajectories(100)
        .max_trajectories(10000)
        .build();

    // -- Initial state: excited state |0> ----------------------------------
    StateVector psi0(2);
    psi0[0] = Scalar{1.0, 0.0};   // |0> component
    psi0[1] = Scalar{0.0, 0.0};   // |1> component

    // -- Run ---------------------------------------------------------------
    auto result = sim.run(psi0, 0.0, T);

    // -- Output -----------------------------------------------------------
    const double exact = 2.0 * std::exp(-gamma * T) - 1.0;
    std::printf("<sigma_z>(T=%.1f) = %.4f +/- %.4f  [exact: %.4f]  N=%zu\n",
        T,
        result.observables[0].mean,
        result.observables[0].sem,
        exact,
        result.total_trajectories);

    return 0;
}
\end{minted}

Expected output (exact value $\approx -0.9003$):

\begin{minted}{bash}
<sigma_z>(T=3.0) = -0.9011 +/- 0.0089  [exact: -0.9003]  N=247
\end{minted}

\begin{notebox}
  The simulation terminates at $N=247$ because the relative standard
  error fell below the 1\% target before reaching the maximum
  trajectory count. The minimum count of 100 was already satisfied.
\end{notebox}

\section{Building the Example}
\label{sec:build-example}

\begin{minted}{bash}
g++ -std=c++17 -O3 -march=native -DNDEBUG \
    -I path/to/LiQuID/include             \
    -L path/to/LiQuID/build               \
    -o two_level_decay two_level_decay.cpp \
    -lliquid
\end{minted}

% -------------------------------------------------------------
\chapter{Layer 0: Foundation}
\label{chap:foundation}
% -------------------------------------------------------------

\section{Type System}
\label{sec:types}

All fundamental types are declared in \filename{liquid/core/types.hpp}
and are available without qualification under \code{namespace liquid}.

\subsection{Scalar Types}
\label{subsec:scalar-types}

\begin{table}[H]
\centering
\caption{Core type aliases. All are defined in \code{liquid/core/types.hpp}.}
\label{tab:core-types}
\begin{tabularx}{\linewidth}{L{22mm}L{38mm}X}
\toprule
\textbf{Alias} & \textbf{Underlying type} & \textbf{Semantic} \\
\midrule
\code{Scalar}  & \code{std::complex<double>} & Complex quantum amplitude \\
\code{Real}    & \code{double}               & Real-valued quantity (time, rate, norm) \\
\code{Dim}     & \code{std::size\_t}         & Hilbert space dimension $N$ \\
\code{Idx}     & \code{std::size\_t}         & Row/column matrix index \\
\code{TrajId}  & \code{std::uint64\_t}       & Unique trajectory identifier \\
\code{Seed}    & \code{std::uint64\_t}       & RNG seed \\
\bottomrule
\end{tabularx}
\end{table}

\begin{invbox}
  \code{Scalar = std::complex<double>} is a \textbf{frozen design decision}.
  The library is not templated on the floating-point type.
  Single-precision (\code{float}) and GPU (\code{half}) support are
  explicitly deferred to a future policy layer and will not appear
  in the public API before v1.0.
\end{invbox}

\subsection{The Imaginary Unit}
\label{subsec:i-unit}

\begin{minted}{cpp}
// In liquid/core/types.hpp:
inline constexpr Scalar i_unit{0.0, 1.0};
\end{minted}

Use \code{i\_unit} in preference to \code{std::complex<double>(0,1)}
or the raw literal \code{1i}. The named constant is self-documenting
and avoids the UDL \code{\_i} suffix's namespace dependency.

\begin{minted}{cpp}
// Correct usage:
using namespace liquid;
Scalar z = 2.5 * i_unit;          // z = 0 + 2.5i
Scalar w = i_unit * H_element;    // w = i * h

// Equivalent but less idiomatic:
Scalar z2 = Scalar{0.0, 2.5};
\end{minted}

\subsection{Enumeration Types}
\label{subsec:enums}

\paragraph{\code{DiagnosticLevel}} Controls the volume of per-trajectory
diagnostic data retained in memory:

\begin{table}[H]
\centering
\caption{\code{DiagnosticLevel} values and their space complexity.}
\label{tab:diag-levels}
\begin{tabularx}{\linewidth}{L{32mm}L{20mm}X}
\toprule
\textbf{Value} & \textbf{Space} & \textbf{Content} \\
\midrule
\code{None}        & $\Order(1)$     & No diagnostics stored; zero per-step overhead \\
\code{Summary}     & $\Order(N_j)$   & Summary statistics + jump event history (default) \\
\code{Full}        & $\Order(N_t)$   & Above + subsampled norm history \\
\code{StepHistory} & $\Order(N_t)$   & All of the above + per-step ODE records (expensive) \\
\bottomrule
\end{tabularx}
\end{table}

\paragraph{\code{TrajectoryStatus}} Trajectory lifecycle states:
\code{Initialized} $\to$ \code{Running} $\to$ \{\code{Completed} $|$ \code{Failed}\}.
The \code{Suspended} state is used by the \code{EnsembleManager} for
adaptive allocation (see \cref{sec:adaptive-allocator}).

\paragraph{\code{SolverType}} \code{RK4Fixed} (Phase~1) and
\code{DormandPrince45} (Phase~3). Select via \code{SimulationBuilder::solver()}.

\section{Configuration Structs}
\label{sec:config}

\subsection{\texttt{PropagatorConfig}}
\label{subsec:propagator-config}

Controls the numerical behaviour of \code{MCWFPropagator}.
All fields have defaults calibrated for quantum optics
problems with $\gamma \sim \omega \sim 1$.

\begin{minted}[linenos]{cpp}
struct PropagatorConfig {
    // ODE stepping
    Real dt_initial      = 1e-3;   // Initial stepsize
    Real dt_min          = 1e-12;  // Floor: triggers Failed status if hit
    Real dt_max          = 1e-1;   // Ceiling: prevents overshooting features

    // Adaptive solver tolerances (used by DormandPrince45)
    Real atol            = 1e-8;   // Absolute local error tolerance
    Real rtol            = 1e-6;   // Relative local error tolerance

    // Jump detection
    Real jump_tol        = 1e-6;   // Accuracy of bisection jump location
    int  jump_bisect_max = 50;     // Max bisection iterations
    int  max_step_rejects = 100;   // Max DOPRI45 rejections before Failed
};
\end{minted}

\begin{warnbox}
  Setting \code{dt\_min} too large will cause the solver to declare
  \code{Failed} for stiff problems. If you observe unexpected failures,
  lower \code{dt\_min} before increasing solver tolerances.
\end{warnbox}

\subsection{\texttt{StoppingCriteria}}
\label{subsec:stopping-criteria}

\begin{minted}[linenos]{cpp}
struct StoppingCriteria {
    Real        target_rel_sem      = 1e-2;   // 1% relative SEM
    std::size_t min_trajectories    = 100;    // Never stop before this
    std::size_t max_trajectories    = 100000; // Hard ceiling
    Real        max_wall_seconds    = 0.0;    // 0 = no limit
    bool        enforce_min         = true;   // Respect min even if converged
};
\end{minted}

The monitor evaluates the following decision hierarchy after each batch:
\begin{enumerate}
  \item If $N < N_{\min}$: \code{Continue} (hard floor).
  \item If $\max_i \mathrm{SEM}_i / |\mu_i| < \varepsilon_{\mathrm{rel}}$:
        \code{Converged}.
  \item If $N \geq N_{\max}$ or elapsed time $> t_{\mathrm{wall}}$:
        \code{BudgetHit}.
  \item If \code{request\_stop()} was called: \code{Forced}.
  \item Otherwise: \code{Continue}.
\end{enumerate}

\section{Random Number Generator}
\label{sec:rng}

\subsection{Algorithm: xoshiro256**}
\label{subsec:xoshiro}

\code{RNGState} implements the \emph{xoshiro256**} algorithm by
Blackman and Vigna (2018). Key properties:

\begin{itemize}
  \item 256-bit state (four \code{uint64\_t}): 32 bytes.
  \item Period $2^{256} - 1 \approx 1.16 \times 10^{77}$.
  \item Passes all BigCrush and PractRand statistical tests.
  \item Approximately $4\times$ faster than \code{std::mt19937\_64}
        with $78\times$ smaller state.
  \item \emph{No heap allocation}. The entire state fits in two
        cache lines.
\end{itemize}

\subsection{Seeding Protocol}
\label{subsec:seeding}

\begin{minted}[linenos]{cpp}
class RNGState {
public:
    // Combine global_seed and traj_id via splitmix64.
    // Different (global_seed, traj_id) pairs are guaranteed to
    // produce statistically independent streams.
    void seed(Seed global_seed, TrajId traj_id) noexcept;

    // Returns a value in (0, 1]  --  note: open at zero.
    // Zero is excluded to prevent spurious immediate jumps at t=0.
    Real draw_uniform() noexcept;

    // Returns an integer in [0, n-1].
    std::size_t draw_int(std::size_t n) noexcept;

    // State serialization for checkpointing.
    void serialize(std::ostream& os) const;
    static RNGState deserialize(std::istream& is);
};
\end{minted}

\begin{invbox}
  \textbf{INV-3 (RNG Exclusivity):} \code{core.rng} is the \emph{sole}
  source of randomness for its trajectory. No global, static, or
  thread-local RNG is accessed anywhere in LiQuID. This guarantees
  that simulation results are byte-identical regardless of thread
  scheduling, provided the seed and trajectory count are the same.
\end{invbox}

The seeding protocol uses \emph{splitmix64} — a bijective hash on
\code{uint64\_t} — to combine the global seed and trajectory identifier:

\begin{equation}
  s_0^{(m)} = \mathrm{splitmix64}(g \oplus m)
  \label{eq:seed}
\end{equation}

where $g$ is the global seed and $m$ is the trajectory index.
Different values of $g$ produce statistically independent ensembles,
enabling reproducible parameter sweeps.

% -------------------------------------------------------------
\chapter{Layer 1: Linear Algebra}
\label{chap:linalg}
% -------------------------------------------------------------

\section{StateVector}
\label{sec:statevector}

\code{StateVector} represents a quantum state $\ket{\psi} \in \complex^N$.
It is stored as a contiguous \code{std::vector<Scalar>} — a flat array
of $N$ complex doubles on the heap.

\subsection{Construction and Lifecycle}
\label{subsec:sv-construction}

\begin{minted}[linenos]{cpp}
// Construct zero-initialized vector of dimension N
StateVector psi(N);

// Construct with explicit fill value
StateVector psi(N, Scalar{0.5, 0.0});

// Copy and move semantics follow standard value semantics.
// Copies are deep (new allocation); moves are O(1).
StateVector phi = psi;                 // deep copy
StateVector chi = std::move(psi);      // O(1) move; psi is now empty
\end{minted}

\subsection{Element Access}
\label{subsec:sv-access}

\begin{minted}{cpp}
psi[0] = Scalar{1.0, 0.0};   // excited state amplitude
psi[1] = Scalar{0.0, 0.0};   // ground state amplitude

// Bounds-checking in Debug builds (assert); unchecked in Release.
// Idx is std::size_t  --  never pass negative integers.
\end{minted}

\subsection{Mathematical Operations}
\label{subsec:sv-math}

\begin{table}[H]
\centering
\caption{\code{StateVector} mathematical interface.}
\label{tab:sv-methods}
\begin{tabularx}{\linewidth}{L{48mm}L{14mm}X}
\toprule
\textbf{Method signature} & \textbf{Alloc?} & \textbf{Computes} \\
\midrule
\code{Real norm\_sq() const}      & No  & $\sum_i |\psi_i|^2$ \\
\code{Real norm() const}          & No  & $\sqrt{\sum_i |\psi_i|^2}$ \\
\code{void normalize()}           & No  & $\ket{\psi} \leftarrow \ket{\psi}/\|\ket{\psi}\|$ \\
\code{void scale(Scalar a)}       & No  & $\ket{\psi} \leftarrow \alpha\ket{\psi}$ \\
\code{void add\_scaled(sv, a)}    & No  & $\ket{\psi} += \alpha\ket{\phi}$ \\
\code{void set\_zero()}           & No  & $\ket{\psi} \leftarrow \mathbf{0}$ \\
\code{void copy\_from(sv)}        & No  & $\ket{\psi} \leftarrow \ket{\phi}$ \\
\code{static Scalar inner(a,b)}   & No  & $\braket{a|b} = \sum_i \bar{a}_i b_i$ \\
\bottomrule
\end{tabularx}
\end{table}

\begin{perfbox}
  \code{add\_scaled()} and \code{copy\_from()} are the two highest-frequency
  operations inside every ODE stepper. They are implemented as tight
  \code{for} loops over the contiguous data array. On x86-64 with
  \code{-O3 -march=native}, the compiler auto-vectorizes them using
  AVX2 (4 complex doubles per SIMD register). Calling \code{apply()}
  (which allocates) inside an ODE loop defeats this optimization.
\end{perfbox}

\subsection{Normalization Invariant}
\label{subsec:sv-norm-invariant}

\begin{invbox}
  \textbf{INV-1 (Normalization):}
  $\|\ket{\psi}\|^2 \approx 1.0$ at $t = t_0$ and immediately after
  every jump. Between jumps, $r < \|\ket{\psi}\|^2 \leq 1.0$
  (strictly decreasing per \cref{eq:norm-decay}).
  \code{normalize()} asserts \code{norm() > 0} in Debug builds and
  is undefined behaviour if called on the zero vector in Release.
\end{invbox}

\section{DenseOperator}
\label{sec:denseoperator}

\code{DenseOperator} represents an $N \times N$ complex matrix
$O \in \complex^{N \times N}$. It is stored \emph{row-major}:
element $(i, j)$ maps to \code{data\_[i * dim\_ + j]}.

\subsection{Storage Layout and Cache Behaviour}
\label{subsec:dense-storage}

Row-major storage is chosen because the dominant operation —
\code{apply\_add(psi, alpha, out)} — sweeps the matrix
row by row, accumulating one output component per row.
For row $i$:

\begin{equation}
  \text{out}[i] \mathrel{+}= \alpha \sum_{j=0}^{N-1}
    \underbrace{A[i \cdot N + j]}_{\text{contiguous in cache}}
    \cdot \psi[j]
  \label{eq:matvec-row}
\end{equation}

The inner sum reads $N$ contiguous \code{Scalar} values from
\code{data\_}, which fits in $N \cdot 16$ bytes. For $N = 64$,
row $i$ occupies 1 kB — exactly one 16-way set of cache lines
on a typical L1 cache.

\subsection{Construction}
\label{subsec:dense-construction}

\begin{minted}[linenos]{cpp}
// Zero-initialized N x N matrix
DenseOperator H(N);

// From flat row-major data (no copy if moved)
std::vector<Scalar> flat_data(N * N, Scalar{0.0, 0.0});
flat_data[0 * N + 0] = Scalar{0.5, 0.0};  // H(0,0)
flat_data[1 * N + 1] = Scalar{-0.5, 0.0}; // H(1,1)
DenseOperator H2(N, std::move(flat_data));

// Factory functions
DenseOperator I  = make_identity(N);          // N x N identity
DenseOperator Z  = make_zero_operator(N);     // N x N zero
DenseOperator LL = adjoint_times(L);          // LdagL  (allocates)
\end{minted}

\subsection{The \texttt{apply\_add} Interface}
\label{subsec:apply-add}

\begin{minted}[linenos]{cpp}
// HOT PATH  --  no allocation, called at every ODE sub-step.
// Computes: out += alpha * A * psi
// Preconditions: psi.size() == dim_, out.size() == dim_
void DenseOperator::apply_add(
    const StateVector& psi,    // input state
    Scalar             alpha,  // scaling factor
    StateVector&       out)    // accumulator (modified in-place)
    const noexcept;

// COLD PATH  --  allocates output vector.
// Use only outside the ODE loop (operator testing, validation).
StateVector DenseOperator::apply(const StateVector& psi) const;
\end{minted}

\begin{warnbox}
  \textbf{Never call \code{apply()} inside an ODE callback.}
  It allocates a new \code{StateVector} on every invocation.
  For $10^4$ trajectories with $10^3$ steps each, this produces
  $10^7$ allocations that defeat the allocator's free-list and
  fragment the heap. Always use \code{apply\_add()} with a
  pre-allocated output vector.
\end{warnbox}

\subsection{Algebraic Operations (Cold Path)}
\label{subsec:dense-algebra}

\begin{minted}[linenos]{cpp}
DenseOperator Adag  = A.adjoint();        // Adag  --  allocates NxN
DenseOperator C     = A.matmul(B);        // A*B  --  allocates NxN
A.add_scaled(B, Scalar{2.0, 0.0});        // A += 2*B  (in-place)

// Free functions
Scalar ev = expectation(O, psi);          // <psi|O|psi> (allocates temp)
\end{minted}

\section{SparseOperator}
\label{sec:sparseoperator}

\code{SparseOperator} stores an $N \times N$ sparse complex matrix in
\emph{Compressed Sparse Row} (CSR) format. It is intended for large
Hilbert spaces where most matrix elements are structurally zero.

\subsection{CSR Format}
\label{subsec:csr-format}

Three arrays define the sparse matrix:

\begin{align}
  \texttt{values}[k]       &\in \complex  &&\text{value of the $k$-th stored element} \\
  \texttt{col\_indices}[k] &\in [0, N)    &&\text{column of the $k$-th stored element} \\
  \texttt{row\_ptr}[i]     &\in [0, \mathrm{nnz}] &&\text{index in \texttt{values} where row $i$ begins}
\end{align}

Row $i$ contains elements at columns
$\texttt{col\_indices}[\texttt{row\_ptr}[i] : \texttt{row\_ptr}[i+1]]$
with values
$\texttt{values}[\texttt{row\_ptr}[i] : \texttt{row\_ptr}[i+1]]$.
The array \texttt{row\_ptr} has length $N+1$; the sentinel
$\texttt{row\_ptr}[N] = \mathrm{nnz}$ gives the total non-zero count.

\subsection{Construction from Triplets (COO $\to$ CSR)}
\label{subsec:triplet-construction}

\begin{minted}[linenos]{cpp}
// Triplet: (row, col, value)
std::vector<Triplet> trips;
trips.push_back({0, 0, Scalar{ 1.0, 0.0}});  // diagonal: A(0,0) = 1
trips.push_back({1, 1, Scalar{-1.0, 0.0}});  // diagonal: A(1,1) = -1
trips.push_back({1, 0, Scalar{ 0.5, 0.0}});  // off-diag: A(1,0) = 0.5

// Duplicate (row,col) entries are automatically summed.
// Elements with |value| < 1e-300 are dropped.
SparseOperator A(N, std::move(trips));

std::printf("nnz = %zu, sparsity = %.4f\n",
    A.nnz(), A.sparsity());  // sparsity = nnz / (N*N)
\end{minted}

\begin{notebox}
  After construction, \code{SparseOperator} is immutable with respect
  to its sparsity pattern. \code{add\_scaled(B, alpha)} internally
  rebuilds the CSR structure from scratch. It is a cold-path operation
  intended for precomputation, not for use inside solvers.
\end{notebox}

\subsection{The Sparse \texttt{apply\_add} Interface}
\label{subsec:sparse-apply-add}

The hot-path interface is identical to \code{DenseOperator}:

\begin{minted}[linenos]{cpp}
// out += alpha * A * psi  (no allocation)
// Inner loop exploits CSR structure: only nnz multiplications
// instead of N^2 for a dense matrix.
void SparseOperator::apply_add(
    const StateVector& psi,
    Scalar             alpha,
    StateVector&       out) const noexcept;
\end{minted}

The complexity of a sparse matrix-vector product is $\Order(\mathrm{nnz})$
rather than $\Order(N^2)$. For the bosonic ladder operator $a$ in an
$N$-dimensional Fock space, $\mathrm{nnz} = N-1$, giving a
$N/(N-1) \approx N$ speedup over the dense case.

\begin{perfbox}
  For Hilbert spaces with $N \leq 64$ and fully connected operators,
  use \code{DenseOperator}: cache reuse dominates and the overhead of
  indirect addressing in CSR is not worth the reduced operation count.
  Switch to \code{SparseOperator} when $\mathrm{nnz} / N^2 \lesssim 0.3$,
  or equivalently when fewer than 30\% of matrix elements are non-zero.
\end{perfbox}

\subsection{Conversion and Free Functions}
\label{subsec:sparse-free}

\begin{minted}[linenos]{cpp}
// Convert dense to sparse (drops elements below zero_tol = 1e-14)
SparseOperator As = to_sparse(A_dense, 1e-14);

// Convert sparse to dense (allocates N*N array)
DenseOperator Ad = A_sparse.to_dense();

// Create sparse identity
SparseOperator I = make_sparse_identity(N);

// Compute AdagA for sparse operator
SparseOperator LdagL = sparse_adjoint_times(L);

// Expectation value <psi|O|psi>
Scalar ev = sparse_expectation(O_sparse, psi);
\end{minted}

% -------------------------------------------------------------
\chapter{Layer 1: ODE Solvers}
\label{chap:ode}
% -------------------------------------------------------------

\section{The Stepper Policy}
\label{sec:stepper-concept}

Both ODE solvers satisfy the \emph{Stepper named requirement} (C++17
\emph{policy class} idiom). Any class satisfying this requirement can
be used as the \code{StepperType} template parameter of
\code{MCWFPropagator}. The requirement is documented below as an
annotated struct; enforcement is via \code{static\_assert} on the
return type of \code{try\_step} inside \code{MCWFPropagator}'s
constructor.

\begin{minted}[linenos]{cpp}
// -- Stepper named requirement (C++17 policy idiom) -----------------------
//
// A type S satisfies Stepper if it provides:
//
//   (1) StepResult try_step(Real t,
//                           const StateVector& psi_in,
//                           StateVector&       psi_out,
//                           Real               dt_suggested,
//                           Real               t_max,
//                           RHSFn&&            rhs_fn);
//
//       Integrates one step from (t, psi_in), writing the result into
//       psi_out.  May shorten the step to respect t_max.
//       Returns a StepResult describing acceptance and the next dt hint.
//       RHSFn is any callable with signature:
//           void(Real t, const StateVector& in, StateVector& accum)
//
//   (2) void reset(Real dt_initial) noexcept;
//
//       Discards all internal stepper state (e.g. FSAL buffers) and
//       resets the working step size.  Must be called after every
//       quantum jump because the state vector is discontinuous.
//
// The library enforces (1) at instantiation time:
//   static_assert(
//       std::is_same_v<
//           decltype(std::declval<S>().try_step(...)),
//           ode::StepResult>,
//       "StepperType must return ode::StepResult from try_step()");

struct StepResult {
    bool accepted;    // true  = step accepted; psi_out is valid
                      // false = step rejected (DOPRI45 only); psi_out undefined
    Real dt_taken;    // actual step length used
    Real dt_next;     // suggested step length for the next call
    Real local_err;   // normalised local error estimate (0.0 for RK4)
};
\end{minted}

The RHS callable passed to \code{try\_step} uses the
\emph{accumulator pattern} — it adds to its output buffer rather
than overwriting it, so the stepper can zero the buffer once and
reuse it across stages:

\begin{minted}{cpp}
// Signature of the RHS callable accepted by every Stepper.
// Computes:  accum += f(t, psi)   -- note: += not =
// The stepper calls set_zero() on accum before the first stage;
// rhs_fn must never zero it internally.
void rhs_fn(Real t, const StateVector& psi, StateVector& accum);
\end{minted}

\section{Fixed-Step RK4 (\texttt{RK4Stepper})}
\label{sec:rk4}

\subsection{Algorithm}
\label{subsec:rk4-algorithm}

The classical fourth-order Runge--Kutta formula integrates
$d\ket{\psi}/dt = f(t, \ket{\psi})$ over one step $h$:

\begin{align}
  k_1 &= f\bigl(t,\; \ket{\psi}\bigr) \label{eq:rk4-k1}\\
  k_2 &= f\bigl(t + h/2,\; \ket{\psi} + (h/2)\,k_1\bigr) \label{eq:rk4-k2}\\
  k_3 &= f\bigl(t + h/2,\; \ket{\psi} + (h/2)\,k_2\bigr) \label{eq:rk4-k3}\\
  k_4 &= f\bigl(t + h,\;\; \ket{\psi} + h\,k_3\bigr) \label{eq:rk4-k4}\\
  \ket{\psi(t+h)} &= \ket{\psi} + \frac{h}{6}\bigl(k_1 + 2k_2 + 2k_3 + k_4\bigr)
  \label{eq:rk4-update}
\end{align}

The local truncation error is $\Order(h^5)$; the global error
accumulates as $\Order(h^4)$.

\subsection{Memory Layout}
\label{subsec:rk4-memory}

\code{RK4Stepper} pre-allocates five \code{StateVector} buffers
($k_1, k_2, k_3, k_4$, and a temporary $\ket{\psi_{\mathrm{tmp}}}$)
on the \emph{first} call to \code{try\_step()}, via
\code{ensure\_scratch(Dim n)}. Subsequent calls reuse them.

\begin{table}[H]
\centering
\caption{\code{RK4Stepper} memory footprint.}
\label{tab:rk4-memory}
\begin{tabular}{lll}
\toprule
\textbf{Buffer} & \textbf{Size} & \textbf{Purpose} \\
\midrule
\code{k1\_}  & $16N$ bytes & Stage 1 derivative \\
\code{k2\_}  & $16N$ bytes & Stage 2 derivative \\
\code{k3\_}  & $16N$ bytes & Stage 3 derivative \\
\code{k4\_}  & $16N$ bytes & Stage 4 derivative \\
\code{tmp\_} & $16N$ bytes & Intermediate state \\
\midrule
Total & $80N$ bytes & (e.g., $80\times64 = 5120$ bytes for $N=64$) \\
\bottomrule
\end{tabular}
\end{table}

\subsection{API}
\label{subsec:rk4-api}

\begin{minted}[linenos]{cpp}
// Construct with fixed stepsize
ode::RK4Stepper stepper(1e-3);    // dt = 1e-3

// Or default-construct and configure via PropagatorConfig
ode::RK4Stepper stepper;
stepper.reset(1e-3);

// StepResult (always accepted for RK4)
struct StepResult {
    bool  accepted;    // true always for RK4
    Real  dt_taken;    // actual dt (may be clamped by t_max)
    Real  dt_next;     // suggestion for next step (same as dt_taken for RK4)
    Real  local_err;   // 0.0 for fixed-step
};
\end{minted}

\section{Adaptive DOPRI45 (\texttt{DormandPrince45})}
\label{sec:dopri45}

\subsection{Algorithm}
\label{subsec:dopri45-algorithm}

The Dormand--Prince 4(5) method is an embedded Runge--Kutta pair
with 7 stage evaluations per accepted step. The Butcher tableau
is (Hairer, N{\o}rsett \& Wanner, Table~5.2):

\begin{equation}
\renewcommand{\arraystretch}{1.3}
\begin{array}{c|ccccccc}
0           &        &           &           &           &           &       &   \\
\tfrac{1}{5}&\tfrac{1}{5} &      &           &           &           &       &   \\
\tfrac{3}{10}&\tfrac{3}{40}&\tfrac{9}{40}& &           &           &       &   \\
\tfrac{4}{5}&\tfrac{44}{45}&-\tfrac{56}{15}&\tfrac{32}{9}& &      &       &   \\
\tfrac{8}{9}&\tfrac{19372}{6561}&-\tfrac{25360}{2187}&\tfrac{64448}{6561}&-\tfrac{212}{729}& & & \\
1           &\tfrac{9017}{3168}&-\tfrac{355}{33}&\tfrac{46732}{5247}&\tfrac{49}{176}&-\tfrac{5103}{18656}& & \\
1           &\tfrac{35}{384}&0&\tfrac{500}{1113}&\tfrac{125}{192}&-\tfrac{2187}{6784}&\tfrac{11}{84}& \\
\hline
\text{5th order} & \tfrac{35}{384} & 0 & \tfrac{500}{1113} & \tfrac{125}{192} & -\tfrac{2187}{6784} & \tfrac{11}{84} & 0 \\
\end{array}
\end{equation}

The \emph{FSAL} (First Same As Last) property means that stage $k_7$ of
an accepted step is identical to stage $k_1$ of the next step.
This reduces the effective cost to \textbf{6 RHS evaluations per accepted
step} rather than 7.

\subsection{PI Stepsize Controller}
\label{subsec:pi-controller}

LiQuID uses the Gustafsson (1991) PI controller for adaptive step
size selection:

\begin{equation}
  h_{n+1} = h_n \cdot S \cdot
    \left(\frac{1}{\varepsilon_n}\right)^{\alpha_{\mathrm{PI}}}
    \cdot
    \left(\varepsilon_{n-1}\right)^{\beta_{\mathrm{PI}}}
  \label{eq:pi-controller}
\end{equation}

where $\varepsilon_n$ is the normalised local error estimate,
$S = 0.9$ is the safety factor, $\alpha_{\mathrm{PI}} = 0.7/5$,
and $\beta_{\mathrm{PI}} = 0.4/5$ (default values calibrated for
the 5th-order DOPRI method). The factor is clipped to $[0.2, 10]$
to prevent runaway step changes.

\subsection{FSAL Implementation}
\label{subsec:fsal}

\begin{minted}[linenos]{cpp}
// In DormandPrince45::try_step()  --  excerpt:

if (/* step accepted */) {
    // Reuse k7 (the 7th stage of this step) as k1 of the next step.
    // k7 = f(t + h, psi5) = f(t_new, psi_new)
    // This is FSAL: saves one RHS evaluation per accepted step.
    k1_.copy_from(k7_);
    fsal_valid_ = true;
    // ...
}

// In compute_stages(), on the first stage:
if (!fsal_valid_) {
    k1_.set_zero();
    rhs_fn(t, y0, k1_);  // Evaluate only if FSAL not available
    ++rhs_count_;
}
// Otherwise, k1_ already holds the correct value  --  no RHS call.
\end{minted}

\begin{perfbox}
  FSAL reduces the typical DOPRI45 cost from 7 to 6 RHS evaluations
  per step (a 14\% saving). The stepper's \code{reset()} method sets
  \code{fsal\_valid\_ = false}, ensuring that after a quantum jump
  (which resets the stepper), the first stage is recomputed correctly
  from the post-jump wavefunction.
\end{perfbox}

\subsection{Configuration}
\label{subsec:dopri45-config}

\begin{minted}[linenos]{cpp}
ode::DormandPrince45Config cfg;
cfg.atol       = 1e-8;    // absolute tolerance
cfg.rtol       = 1e-6;    // relative tolerance
cfg.dt_initial = 1e-3;
cfg.dt_min     = 1e-12;   // floor
cfg.dt_max     = 1e-1;    // ceiling
cfg.safety     = 0.9;     // PI controller safety factor
cfg.alpha      = 0.7/5.0; // PI integral gain
cfg.beta       = 0.4/5.0; // PI derivative gain
cfg.max_rejects = 100;    // max rejections before failure

ode::DormandPrince45 stepper(cfg);
\end{minted}

% -------------------------------------------------------------
\chapter{Layer 2: Quantum System}
\label{chap:system}
% -------------------------------------------------------------

\section{LindbladSet}
\label{sec:lindbladset}

\code{LindbladSet<Tag>} stores the collection of jump operators
$\{L_k\}_{k=1}^{K}$ and their precomputed aggregate:

\begin{equation}
  \Gamma = \sum_{k=1}^{K} L_k^\dagger L_k
  \label{eq:gamma-aggregate}
\end{equation}

Two specializations exist:
\begin{itemize}
  \item \code{LindbladSet<DenseTag>} — uses \code{DenseOperator}.
  \item \code{LindbladSet<SparseTag>} — uses \code{SparseOperator}.
\end{itemize}

Both have an identical public interface.

\subsection{Why Precompute $\Gamma$?}
\label{subsec:gamma-precompute}

The non-Hermitian Hamiltonian is $\Heff = H - \tfrac{i}{2}\Gamma$.
At every ODE sub-step, the propagator must evaluate:

\begin{equation}
  -i\Heff\ket{\psi}
    = -iH\ket{\psi} - \frac{1}{2}\Gamma\ket{\psi}
  \label{eq:heff-apply}
\end{equation}

Without precomputation of $\Gamma$, this would require $K$ separate
matrix-vector products $L_k^\dagger L_k \ket{\psi}$ at every step —
$\Order(K \cdot N^2)$ work. Precomputing $\Gamma$ reduces this to a
single product $\Gamma\ket{\psi}$ — $\Order(N^2)$, independent of $K$.
For $K = 5$ operators, this is a $5\times$ reduction in the dominant
inner-loop cost.

\subsection{Construction and Immutability}
\label{subsec:lindblad-construction}

\begin{minted}[linenos]{cpp}
// Dense specialization
std::vector<DenseOperator> ops;
ops.push_back(make_sigma_minus(gamma_1));
ops.push_back(make_sigma_minus(gamma_2));  // two decay channels

LindbladSet<DenseTag> L(std::move(ops));
// ^ Internally calls precompute_gamma():
//     gamma_ = L[0]dag*L[0] + L[1]dag*L[1]

// Copy is disabled (prevents accidental stale-gamma bugs):
// LindbladSet<DenseTag> L2 = L;  // COMPILE ERROR

// Move is allowed:
LindbladSet<DenseTag> L3 = std::move(L);  // OK
\end{minted}

\begin{invbox}
  \textbf{INV-$\Gamma$:}
  $\Gamma \equiv \sum_k L_k^\dagger L_k$ is established at construction
  and is immutable thereafter. The copy constructor is deleted to
  prevent the creation of a \code{LindbladSet} with a stale
  $\Gamma$. Parameter sweeps must construct new \code{LindbladSet}
  instances for each parameter value.
\end{invbox}

\subsection{Hot-Path Interface}
\label{subsec:lindblad-hot}

\begin{minted}[linenos]{cpp}
// Called at every ODE sub-step.
// Computes: out += -(i/2) * Gamma * psi
// where Gamma = sum_k LkdagLk (precomputed)
void LindbladSet<DenseTag>::apply_decay(
    const StateVector& psi,
    StateVector& out) const noexcept;
// alpha = -(i/2) = Scalar{0.0, -0.5} is hard-coded
\end{minted}

\subsection{Cold-Path Interface}
\label{subsec:lindblad-cold}

\begin{minted}[linenos]{cpp}
std::size_t K = L.num_channels();      // K = number of operators
Dim         N = L.hilbert_dim();       // Hilbert space dimension

// Evaluate L[k] * psi, store in out (used for jump probability)
StateVector scratch(N);
L.apply_channel(k, psi, scratch);      // out = L[k] * psi

// Compute p[k] = ||L[k] * psi||^2 for all k
std::vector<Real> probs(K);
Real total = L.all_probabilities(psi, scratch, probs);
// probs[k] = ||L[k] * psi||^2, total = sum(probs)
\end{minted}

\section{OpenSystem}
\label{sec:opensystem}

\code{OpenSystem<TimeTag, StorageTag>} encapsulates the complete
physics description of the open quantum system.

\subsection{Dense Time-Independent System}
\label{subsec:opensystem-dense-ti}

This is the primary type for Phase~1 simulations and is aliased
as \code{DenseOpenSystem}:

\begin{minted}[linenos]{cpp}
// Construct from Hamiltonian and Lindblad operators.
// Preconditions:
//   H.size() == L.hilbert_dim()
//   Both must be constructed before being moved in.
DenseOpenSystem sys(std::move(H), std::move(L));
// ^ Immediately calls precompute_H_eff():
//     H_eff_ = H - (i/2) * Gamma

// HOT PATH: d|psi>/dt = -i * H_eff * psi
// Sets out = -i * H_eff * psi (out is overwritten, not accumulated)
sys.apply_Heff(psi, out);

// COLD PATH: jump operations
std::vector<Real> probs(sys.num_channels());
sys.jump_probabilities(psi, scratch, probs);  // fills probs
sys.apply_jump(k, psi, scratch);              // psi <- Lk*psi/||Lk*psi||

// Introspection
sys.hilbert_dim();     // N
sys.num_channels();    // K
sys.H_eff();           // const ref to precomputed H_eff
\end{minted}

\subsection{Precomputed \texorpdfstring{$\Heff$}{Heff}}
\label{subsec:heff-precompute}

\begin{minted}[linenos]{cpp}
// Excerpt from precompute_H_eff():
void precompute_H_eff() {
    H_eff_ = H_;                            // H_eff = H
    constexpr Scalar neg_i_half{0.0, -0.5}; // -(i/2)
    H_eff_.add_scaled(L_.decay_operator(),  // += -(i/2) * Gamma
                      neg_i_half);
}
// Result: H_eff = H - (i/2) * Gamma  (frozen at construction)
\end{minted}

The constant $-i/2 = \{0.0, -0.5\}$ appears in three places:
\code{precompute\_H\_eff()}, \code{apply\_decay()}, and the
\code{apply\_Heff()} call to \code{H\_eff\_.apply\_add(psi, -i\_unit, out)}.

\subsection{Sparse Time-Independent System}
\label{subsec:opensystem-sparse}

\begin{minted}[linenos]{cpp}
// Construct with sparse operators
SparseOpenSystem sys_sparse(
    SparseOperator H_sparse,
    LindbladSet<SparseTag> L_sparse);

// Interface identical to DenseOpenSystem:
sys_sparse.apply_Heff(psi, out);
sys_sparse.apply_jump(k, psi, scratch);
\end{minted}

\subsection{Time-Dependent Hamiltonian}
\label{subsec:opensystem-timedep}

For $H(t) = H_{\mathrm{static}} + \sum_j f_j(t)\,H_{\mathrm{drive},j}$:

\begin{minted}[linenos]{cpp}
// DriveTerm: one time-dependent drive term
DriveTerm<DenseTag> drive;
drive.envelope = [](Real t) { return std::cos(omega_d * t); };
drive.H_drive  = make_H_drive_dense(N, coupling_g);

// Construct time-dependent open system
auto sys_td = OpenSystem<TimeDependentTag, DenseTag>(
    std::move(H_static),
    std::move(L),
    std::vector<DriveTerm<DenseTag>>{std::move(drive)});

// HOT PATH (now takes time argument):
// out = -i * [H_static + sum_j f_j(t)*H_drive_j - (i/2)*Gamma] * psi
sys_td.apply_Heff(t, psi, out);
\end{minted}

% -------------------------------------------------------------
\chapter{Layer 3: Trajectory}
\label{chap:trajectory}
% -------------------------------------------------------------

\section{TrajectoryState}
\label{sec:trajectory-state}

\code{TrajectoryState} is the complete mutable record of a single
MCWF trajectory. It is always created via the factory function
\code{make\_trajectory\_state()} and modified exclusively by
\code{MCWFPropagator::advance()}.

\subsection{CoreState}
\label{subsec:corestate}

\begin{minted}[linenos]{cpp}
struct CoreState {
    StateVector      psi;       // current wavefunction  [INV-1]
    Real             t;         // current time          [INV-4]
    Real             t_final;   // end time  --  immutable  [INV-7]
    Real             r;         // jump threshold in (0,1]  [INV-2]
    Real             dt_last;   // warm-start hint for next step
    Real             err_last;  // last local error estimate
    RNGState         rng;       // trajectory-exclusive RNG [INV-3]
    TrajId           traj_id;   // immutable identifier  [INV-6]
    TrajectoryStatus status;    // lifecycle enum
};
\end{minted}

The seven invariants of \code{CoreState} are defined in
\cref{sec:types} and enforced exclusively by \code{MCWFPropagator}.

\subsection{Factory Function}
\label{subsec:trajectory-factory}

\begin{minted}[linenos]{cpp}
// The ONLY legal way to create a TrajectoryState.
TrajectoryState ts = make_trajectory_state(
    psi0,                           // initial state (normalized)
    0.0,                            // t0
    10.0,                           // t_final
    traj_id,                        // unique ID (e.g., 0, 1, 2, ...)
    global_seed,                    // seed for this trajectory's RNG
    DiagnosticLevel::Summary        // optional, default Summary
);
// Postconditions:
//   ts.core.psi.norm_sq() ~= 1.0
//   ts.core.r is a fresh draw in (0,1]
//   ts.core.status == TrajectoryStatus::Initialized
\end{minted}

\subsection{DiagnosticRecord}
\label{subsec:diagnostic-record}

\begin{minted}[linenos]{cpp}
if (ts.has_diagnostics()) {
    const auto& diag = ts.diag();

    std::printf("Total steps : %llu\n", diag.total_steps);
    std::printf("Total jumps : %llu\n", diag.total_jumps);
    std::printf("Mean dt     : %.3e\n", diag.mean_dt);
    std::printf("Min dt      : %.3e\n", diag.min_dt);
    std::printf("Wall time   : %.3f s\n", diag.wall_time_s);

    // Jump events (requires DiagnosticLevel::Summary or higher)
    for (const auto& jump : diag.jumps) {
        std::printf("  Jump at t=%.4f  channel=%zu  norm_before=%.6f\n",
            jump.t, jump.channel, jump.norm_before);
    }
}
\end{minted}

\section{MCWFPropagator}
\label{sec:mcwf-propagator}

\code{MCWFPropagator<SystemType, StepperType>} implements the
MCWF algorithm. It is the computational engine of LiQuID.

\subsection{Per-Instance Scratch Space}
\label{subsec:propagator-scratch}

Five \code{StateVector} buffers are allocated at construction time
and reused for every call to \code{advance()}:

\begin{table}[H]
\centering
\caption{\code{MCWFPropagator} per-instance scratch buffers.}
\label{tab:propagator-scratch}
\begin{tabular}{lll}
\toprule
\textbf{Member} & \textbf{Size} & \textbf{Purpose} \\
\midrule
\code{psi\_scratch\_}   & $16N$ bytes & General-purpose temporary \\
\code{psi\_proposed\_}  & $16N$ bytes & ODE step output before jump check \\
\code{psi\_before\_}    & $16N$ bytes & State saved before each ODE step \\
\code{accum\_tmp\_}     & $16N$ bytes & RHS accumulation temporary \\
\code{prob\_scratch\_}  & $8K$ bytes  & Per-channel jump probabilities \\
\midrule
Total & $64N + 8K$ bytes & (e.g., $\sim 4.1\,\mathrm{kB}$ for $N=64, K=3$) \\
\bottomrule
\end{tabular}
\end{table}

\begin{invbox}
  \code{MCWFPropagator} is \textbf{not thread-safe}. Each thread must
  own its own propagator instance. Multiple propagators may safely
  share a single \code{const SystemType*} pointer because
  \code{OpenSystem} is immutable and \code{apply\_Heff()} is
  declared \code{const noexcept}.
\end{invbox}

\subsection{Direct Usage}
\label{subsec:propagator-direct}

Direct use of \code{MCWFPropagator} is appropriate when fine-grained
control over a single trajectory is required — for example, in
debugging, visualization, or single-trajectory benchmarking:

\begin{minted}[linenos]{cpp}
#include "liquid/liquid.hpp"
using namespace liquid;
using namespace liquid::ode;
using namespace liquid::trajectory;

int main() {
    // Build system
    DenseOperator H(2);
    H(0,0) = Scalar{0.5, 0.0}; H(1,1) = Scalar{-0.5, 0.0};
    DenseOperator L(2);
    L(1,0) = Scalar{1.0, 0.0};

    DenseOpenSystem sys(H, LindbladSet<DenseTag>({L}));

    // Create trajectory
    StateVector psi0(2);
    psi0[0] = Scalar{1.0, 0.0};

    auto ts = make_trajectory_state(psi0, 0.0, 5.0,
                                    /*traj_id=*/0,
                                    /*seed=*/42,
                                    DiagnosticLevel::Full);

    // Create propagator with RK4 stepper
    PropagatorConfig cfg;
    cfg.dt_initial = 1e-3;
    MCWFPropagator<DenseOpenSystem, RK4Stepper> prop(
        &sys,
        RK4Stepper(cfg.dt_initial),
        cfg);

    // Event-driven loop: inspect state after every jump
    using liquid::trajectory::StepOutcome;
    StepOutcome outcome = StepOutcome::Advanced;

    while (outcome != StepOutcome::Completed &&
           outcome != StepOutcome::Failed) {
        outcome = prop.advance(ts, ts.core.t_final);

        if (outcome == StepOutcome::JumpOccurred) {
            std::printf("Jump at t=%.4f  ||psi||^2 = %.6f\n",
                ts.core.t,
                ts.core.psi.norm_sq());
        }
    }

    if (ts.is_completed()) {
        std::printf("Trajectory completed. Total jumps: %llu\n",
            ts.diag().total_jumps);
    } else {
        std::printf("Trajectory FAILED at t=%.4f\n", ts.core.t);
    }

    return 0;
}
\end{minted}

\subsection{Convenience: \texttt{run\_to\_completion}}
\label{subsec:run-to-completion}

\begin{minted}[linenos]{cpp}
// Equivalent to calling advance() in a loop until Completed or Failed.
// Records total wall time in ts.diag().wall_time_s.
TrajectoryStatus status = prop.run_to_completion(ts);

if (status == TrajectoryStatus::Completed) {
    // Compute observable
    DenseOperator sz(2);
    sz(0,0) = Scalar{1,0}; sz(1,1) = Scalar{-1,0};
    Scalar ev = expectation(sz, ts.core.psi);
    std::printf("<sz> = %.4f\n", ev.real());
}
\end{minted}

% -------------------------------------------------------------
\chapter{Layer 4: Ensemble Management}
\label{chap:ensemble}
% -------------------------------------------------------------

\section{ObservableAccumulator}
\label{sec:observable-accumulator}

\code{ObservableAccumulator} bridges the physics layer (wavefunctions)
and the statistics layer (Welford online estimators).

\subsection{Observable Definitions}
\label{subsec:observable-defs}

An observable is defined as a callable \code{ObservableFn}:

\begin{minted}[linenos]{cpp}
using ObservableFn = std::function<Real(const StateVector&)>;

// Example: expectation value of sigma_z
DenseOperator sz(2);
sz(0,0) = Scalar{1,0}; sz(1,1) = Scalar{-1,0};

// The ObservableFn must normalize internally:
//   <O> = <psi|O|psi> / <psi|psi>
ObservableDef sz_def;
sz_def.name = "sigma_z";
sz_def.eval = [op = sz](const StateVector& psi) -> Real {
    const Real nrm2 = psi.norm_sq();
    if (nrm2 < 1e-14) return 0.0;
    return expectation(op, psi).real() / nrm2;
};
\end{minted}

\begin{notebox}
  The \code{SimulationBuilder::observe("name", DenseOperator)} convenience
  method automatically wraps the operator in a normalized expectation
  value callable. Use the raw \code{ObservableFn} interface only when
  computing custom quantities (e.g., photon number variance,
  Wigner function samples, or operator correlators).
\end{notebox}

\section{ConvergenceMonitor}
\label{sec:convergence-monitor}

\code{ConvergenceMonitor} answers the question ``should the ensemble stop?''
It is stateless except for an atomic stop flag:

\begin{minted}[linenos]{cpp}
StoppingCriteria sc;
sc.target_rel_sem   = 0.01;    // 1% relative SEM
sc.min_trajectories = 100;
sc.max_trajectories = 50000;

ConvergenceMonitor monitor(sc);

// After each trajectory batch:
ConvergenceReport report = monitor.evaluate(accum, elapsed_seconds);

switch (report.decision) {
    case StoppingDecision::Continue:   /* keep going */ break;
    case StoppingDecision::Converged:
        std::printf("Converged at N=%zu, rel_sem=%.4f\n",
            report.trajectories_completed, report.current_rel_sem);
        break;
    case StoppingDecision::BudgetHit:
        std::printf("Budget exhausted: %s\n",
            report.reason.c_str());
        break;
    case StoppingDecision::Forced:
        std::printf("Stop requested externally.\n");
        break;
}
\end{minted}

\section{Adaptive Allocator and Policies}
\label{sec:adaptive-allocator}

The allocator policy is a \code{std::function<AllocationDecision(EnsembleSummary)>}.
The default is \code{make\_uniform\_policy()}: always spawn a new trajectory.

\begin{minted}[linenos]{cpp}
// Uniform policy (default): always new trajectory, round-robin seeds
AllocatorPolicy uniform = make_uniform_policy(global_seed);

// Greedy policy: prioritise seeds that minimise current SEM
AllocatorPolicy greedy = make_greedy_policy(global_seed);

// ML policy (Phase 6): 5-feature MLP guiding seed selection
AllocatorPolicy ml = ensemble::ml::make_ml_policy(global_seed);
\end{minted}

The \code{TrajectorySummary} struct (with fields \code{total\_jumps},
\code{total\_steps}, \code{wall\_time\_seconds}, \code{mean\_dt})
is the primary feature vector for the ML policy. It is also the
data format for logging and post-hoc analysis.

\section{EnsembleManager}
\label{sec:ensemble-manager}

\code{EnsembleManager<SystemType, StepperType>} orchestrates
multi-threaded trajectory dispatch via OpenMP:

\begin{minted}[linenos]{cpp}
// Typically not instantiated directly: use SimulationBuilder.
// Shown here for illustration of the threading model.
EnsembleManager<DenseOpenSystem, RK4Stepper> manager(
    &sys,                            // shared, read-only
    std::move(observable_defs),
    stopping_criteria,
    ensemble_config,                 // num_threads, seed, diag_level
    make_uniform_policy(42));

EnsembleResult result = manager.run(psi0, 0.0, t_final);

std::printf("N=%zu  failed=%zu  jumps_mean=%.2f  wall=%.2fs\n",
    result.total_trajectories,
    result.failed_trajectories,
    result.mean_jumps_per_trajectory,
    result.total_wall_time_seconds);
\end{minted}

% -------------------------------------------------------------
\chapter{Layer 5: Simulation Interface}
\label{chap:simulation}
% -------------------------------------------------------------

\section{SimulationBuilder}
\label{sec:simulation-builder}

\code{SimulationBuilder} is the primary user-facing entry point.
It uses the fluent builder pattern to accumulate configuration
and produce a type-erased \code{Simulation} handle via \code{build()}.

\subsection{Complete Builder Reference}
\label{subsec:builder-reference}

\begin{minted}[linenos]{cpp}
Simulation sim = SimulationBuilder{}
    // -- Physics ------------------------------------------------------
    .hamiltonian(H)                  // DenseOperator: system Hamiltonian
    .collapse_operator(L1)           // Add first jump operator
    .collapse_operator(L2)           // Chain multiple operators
    // -- Observables --------------------------------------------------
    .observe("sz", sz)               // Dense operator observable
    .observe("n_photon", [](const StateVector& psi) -> Real {
        // Custom observable (e.g., photon number)
        Real n = 0.0;
        for (Idx i = 0; i < psi.size(); ++i)
            n += static_cast<Real>(i) * std::norm(psi[i]);
        return n / psi.norm_sq();
    })
    // -- ODE solver ---------------------------------------------------
    .dt(1e-3)                        // PropagatorConfig::dt_initial
    .solver(SolverType::DormandPrince45)  // default: RK4Fixed
    .atol(1e-8)                      // DOPRI45: absolute tolerance
    .rtol(1e-6)                      // DOPRI45: relative tolerance
    // -- Stopping -----------------------------------------------------
    .stop_at_sem(0.01)               // StoppingCriteria::target_rel_sem
    .min_trajectories(100)
    .max_trajectories(10000)
    .max_wall_time(60.0)             // seconds; 0 = no limit
    // -- Parallelism --------------------------------------------------
    .num_threads(4)                  // requires -fopenmp
    .seed(42)
    // -- Diagnostics --------------------------------------------------
    .diagnostic_level(DiagnosticLevel::Summary)
    .build();
\end{minted}

\subsection{Running and Inspecting Results}
\label{subsec:running-results}

\begin{minted}[linenos]{cpp}
auto result = sim.run(psi0, t0, t_final);

// Per-observable statistics
for (const auto& obs : result.observables) {
    std::printf("%-20s  mean=%+.6f  sem=%.2e  rel_sem=%.2e  N=%zu\n",
        obs.name.c_str(),
        obs.mean,
        obs.sem,
        obs.rel_sem,
        result.total_trajectories);
}

// Ensemble metadata
std::printf("Convergence: %s\n", result.convergence.reason.c_str());
std::printf("Threads used: %zu\n", result.num_threads_used);
std::printf("Mean jumps/traj: %.2f\n", result.mean_jumps_per_trajectory);
\end{minted}

\section{Parameter Sweep}
\label{sec:parameter-sweep}

\code{ParameterSweepBuilder} runs the same simulation at multiple
parameter values and collects results with full statistical metadata.

\begin{minted}[linenos]{cpp}
#include "liquid/simulation/parameter_sweep.hpp"
using namespace liquid;

int main() {
    // Shared operators (not parameter-dependent)
    DenseOperator sz(2);
    sz(0,0) = Scalar{1,0}; sz(1,1) = Scalar{-1,0};

    StateVector psi0(2);
    psi0[0] = Scalar{1.0, 0.0};  // excited state

    // Define the sweep
    auto sweep = ParameterSweepBuilder{}
        .parameter("gamma", {0.1, 0.5, 1.0, 2.0, 5.0, 10.0})
        .simulation_factory([&sz](double gamma) {
            DenseOperator H(2);
            H(0,0) = Scalar{0.5, 0.0}; H(1,1) = Scalar{-0.5, 0.0};

            DenseOperator L(2);
            L(1,0) = Scalar{std::sqrt(gamma), 0.0};

            return SimulationBuilder{}
                .hamiltonian(std::move(H))
                .collapse_operator(std::move(L))
                .observe("sigma_z", sz)
                .seed(42)
                .dt(1e-3)
                .stop_at_sem(0.01)
                .min_trajectories(100)
                .max_trajectories(5000)
                .build();
        })
        .initial_state(psi0)
        .time_interval(0.0, 5.0)
        .build();

    SweepResult result = sweep.run();

    // Save to CSV: param,obs_name,mean,sem,rel_sem,N_trajectories
    result.save_csv("sweep_gamma.csv");

    // Access individual points
    for (const auto& pt : result.points) {
        const auto& obs = pt.result.observables[0];
        std::printf("gamma=%.2f  <sz>=%.4f +/- %.4f  (N=%zu)\n",
            pt.param_value, obs.mean, obs.sem,
            pt.result.total_trajectories);
    }

    return 0;
}
\end{minted}

% ============================================================
%  PART III — EXAMPLES AND VALIDATION
% ============================================================
\part{Examples and Validation}

% -------------------------------------------------------------
\chapter{Physics Examples}
\label{chap:examples}
% -------------------------------------------------------------

\section{Example 1: Two-Level Spontaneous Decay}
\label{sec:ex-two-level}

The paradigmatic open quantum system: a two-level atom at frequency
$\omega$, decaying at rate $\gamma$ via spontaneous emission.

\textbf{Exact solution:}
\begin{equation}
  \langle\sigma_z(t)\rangle = (P_e(0) - P_g(0))\,e^{-\gamma t}
    + (P_g(0) - P_e(0))
  \label{eq:two-level-exact-general}
\end{equation}

For initial excited state $\ket{e}$: $P_e(0) = 1$, so
$\langle\sigma_z\rangle = 2e^{-\gamma t} - 1$.

\begin{minted}[linenos]{cpp}
#include "liquid/liquid.hpp"
#include <cmath>
#include <cstdio>
using namespace liquid;
using namespace liquid::ensemble;

int main() {
    const double omega = 1.0;
    const double gamma = 1.0;
    const double T     = 5.0;

    // Hamiltonian
    DenseOperator H(2);
    H(0,0) = Scalar{ omega/2, 0}; H(1,1) = Scalar{-omega/2, 0};

    // Jump operator: sigma_minus = sqrt(gamma) |g><e|
    DenseOperator L(2);
    L(1,0) = Scalar{std::sqrt(gamma), 0};

    // Observable: sigma_z = |e><e| - |g><g|
    DenseOperator sz(2);
    sz(0,0) = Scalar{1, 0}; sz(1,1) = Scalar{-1, 0};

    // Build and run
    Simulation sim = SimulationBuilder{}
        .hamiltonian(std::move(H))
        .collapse_operator(std::move(L))
        .observe("sigma_z", std::move(sz))
        .seed(1337)
        .dt(1e-3)
        .stop_at_sem(0.005)          // 0.5% SEM target
        .min_trajectories(200)
        .max_trajectories(20000)
        .build();

    StateVector psi0(2);
    psi0[0] = Scalar{1.0, 0.0};  // excited state

    auto result = sim.run(psi0, 0.0, T);

    const double exact = 2.0 * std::exp(-gamma * T) - 1.0;
    const double error = std::abs(result.observables[0].mean - exact);
    const double nsig  = error / result.observables[0].sem;

    std::printf("=== Two-Level Spontaneous Decay ===\n");
    std::printf("  <sigma_z>  MCWF : %+.6f +/- %.6f\n",
        result.observables[0].mean,
        result.observables[0].sem);
    std::printf("  <sigma_z>  Exact: %+.6f\n", exact);
    std::printf("  Deviation        : %.2f sigma\n", nsig);
    std::printf("  Trajectories     : %zu\n", result.total_trajectories);
    std::printf("  Mean jumps/traj  : %.2f\n",
        result.mean_jumps_per_trajectory);

    // Validation: result must be within 3 sigma of exact
    return (nsig < 3.0) ? 0 : 1;
}
\end{minted}

\section{Example 2: Sparse Jaynes--Cummings Steady State}
\label{sec:ex-jc}

The Jaynes--Cummings model describes a two-level atom coupled to a
single electromagnetic cavity mode. At steady state with cavity decay
$\kappa$ and atomic decay $\gamma$, the mean photon number
$\langle a^\dagger a \rangle_{\mathrm{ss}}$ can be computed analytically.

\textbf{System Hamiltonian} (rotating frame):
\begin{equation}
  H = \Delta_c\,a^\dagger a + \Delta_a\,\sigma_+\sigma_-
    + g\,(a\,\sigma_+ + a^\dagger\,\sigma_-)
    + \Omega\,(a + a^\dagger)
  \label{eq:jc-hamiltonian}
\end{equation}

where $g$ is the vacuum Rabi coupling, $\Omega$ the drive amplitude,
and $\Delta_{c,a}$ the cavity and atom detunings.


% -----------------------------------------------------------------------------
\section{Example 2: Driven Two-Level System}
\label{sec:ex-driven}
% -----------------------------------------------------------------------------

A coherent drive (Rabi frequency $\Omega$) competes with spontaneous emission
(rate $\gamma$). In the rotating frame at the resonance frequency, the
Hamiltonian is $H = \Omega\,\sigma_x$ and the single jump operator is
$L = \sqrt{\gamma}\,\sigma_-$. The optical Bloch equations give the
steady-state excited population:

\begin{equation}
  \varrho_{ee}^{\mathrm{ss}}
    = \frac{4\Omega^2}{\gamma^2 + 8\Omega^2}
  \label{eq:driven-ss}
\end{equation}

Note the factor of 4 (not 1): it arises because
$H = \Omega\,\sigma_x$ (not $\tfrac{\Omega}{2}\,\sigma_x$).
\cref{eq:driven-ss} is verified numerically for two drive regimes.

\begin{minted}[linenos]{cpp}
#include "liquid/liquid.hpp"
#include <cmath>
#include <cstdio>
using namespace liquid;
using namespace liquid::ensemble;

// Analytic steady-state excited population for H = Omega*sigma_x,
// L = sqrt(gamma)*sigma_minus.
static double analytic_rho_ee_ss(double Omega, double gamma) {
    const double r = Omega / gamma;
    return 4.0 * r * r / (1.0 + 8.0 * r * r);
}

// Build the rotating-frame driven open system
static DenseOpenSystem make_driven_system(double Omega, double gamma) {
    // H = Omega * sigma_x  (resonant drive, rotating frame)
    DenseOperator H(2);
    H(0, 1) = Scalar{Omega, 0.0};
    H(1, 0) = Scalar{Omega, 0.0};

    // L = sqrt(gamma) * sigma_minus
    DenseOperator L(2);
    L(1, 0) = Scalar{std::sqrt(gamma), 0.0};

    std::vector<DenseOperator> ops;
    ops.push_back(std::move(L));
    return DenseOpenSystem(std::move(H), LindbladSet<DenseTag>(std::move(ops)));
}

int main() {
    const double gamma  = 1.0;
    const double T_ss   = 20.0;  // long enough to reach steady state
    const int    N_traj = 3000;

    // -- Regime 1: weak drive (Omega = 0.1*gamma) --------------------------
    {
        const double Omega     = 0.1 * gamma;
        const double exact_ss  = analytic_rho_ee_ss(Omega, gamma);

        auto sys = make_driven_system(Omega, gamma);

        PropagatorConfig cfg;
        cfg.dt_initial = 5e-4;  // small dt: Rabi period = pi/(Omega) ~ 31

        double sum_rho = 0.0, sum_rho2 = 0.0;
        StateVector psi0(2);
        psi0[0] = Scalar{1.0, 0.0};  // start excited

        for (int traj = 0; traj < N_traj; ++traj) {
            liquid::ode::RK4Stepper stepper(cfg.dt_initial);
            liquid::trajectory::MCWFPropagator<
                DenseOpenSystem, liquid::ode::RK4Stepper> prop(&sys, stepper, cfg);

            auto ts = liquid::trajectory::make_trajectory_state(
                psi0, 0.0, T_ss,
                static_cast<TrajId>(traj), 54321ULL,
                DiagnosticLevel::None);

            prop.run_to_completion(ts);

            const Real ns      = ts.core.psi.norm_sq();
            const double rho_e = ns > 1e-30
                ? std::norm(ts.core.psi[0]) / ns
                : 0.0;
            sum_rho  += rho_e;
            sum_rho2 += rho_e * rho_e;
        }

        const double mean = sum_rho / N_traj;
        const double var  = sum_rho2 / N_traj - mean * mean;
        const double sem  = std::sqrt(std::max(0.0, var) / N_traj);

        std::printf("Weak drive (Omega/gamma = %.2f):\n", Omega / gamma);
        std::printf("  rho_ee(MCWF)  = %.4f +/- %.4f\n", mean, sem);
        std::printf("  rho_ee(exact) = %.4f\n", exact_ss);
        std::printf("  |error|       = %.4f  (%.1f sigma)\n\n",
            std::abs(mean - exact_ss),
            std::abs(mean - exact_ss) / sem);
    }

    // -- Regime 2: strong drive (Omega = gamma) ----------------------------
    {
        const double Omega     = 1.0 * gamma;
        const double exact_ss  = analytic_rho_ee_ss(Omega, gamma);

        auto sys = make_driven_system(Omega, gamma);

        PropagatorConfig cfg;
        cfg.dt_initial = 5e-4;

        double sum_rho = 0.0, sum_rho2 = 0.0;
        StateVector psi0(2);
        psi0[0] = Scalar{1.0, 0.0};

        for (int traj = 0; traj < N_traj; ++traj) {
            liquid::ode::RK4Stepper stepper(cfg.dt_initial);
            liquid::trajectory::MCWFPropagator<
                DenseOpenSystem, liquid::ode::RK4Stepper> prop(&sys, stepper, cfg);

            auto ts = liquid::trajectory::make_trajectory_state(
                psi0, 0.0, T_ss,
                static_cast<TrajId>(traj), 99999ULL,
                DiagnosticLevel::None);

            prop.run_to_completion(ts);

            const Real ns      = ts.core.psi.norm_sq();
            const double rho_e = ns > 1e-30
                ? std::norm(ts.core.psi[0]) / ns
                : 0.0;
            sum_rho  += rho_e;
            sum_rho2 += rho_e * rho_e;
        }

        const double mean = sum_rho / N_traj;
        const double var  = sum_rho2 / N_traj - mean * mean;
        const double sem  = std::sqrt(std::max(0.0, var) / N_traj);

        std::printf("Strong drive (Omega/gamma = %.2f):\n", Omega / gamma);
        std::printf("  rho_ee(MCWF)  = %.4f +/- %.4f\n", mean, sem);
        std::printf("  rho_ee(exact) = %.4f  [saturation: 1/3 limit]\n", exact_ss);
        std::printf("  |error|       = %.4f  (%.1f sigma)\n",
            std::abs(mean - exact_ss),
            std::abs(mean - exact_ss) / sem);
    }

    return 0;
}
\end{minted}

\begin{notebox}
  \textbf{Rabi oscillations limit.} Setting $\gamma \to 0$ with $\Omega > 0$
  yields nearly closed Rabi oscillations: $\langle\sigma_z(t)\rangle
  = \cos(2\Omega t)\,e^{-\gamma t/2}$.
  \code{LindbladSet} requires at least one operator, so use a very small
  $\gamma \approx 10^{-2}$ to approximate the closed-system limit numerically.
\end{notebox}

% -----------------------------------------------------------------------------
\section{Example 3: Cavity QED --- Jaynes--Cummings Model}
\label{sec:ex-cavity-qed}
% -----------------------------------------------------------------------------

The Jaynes--Cummings (JC) model describes a single two-level atom
coupled to one cavity mode. The composite Hilbert space is
$\mathcal{H} = \mathcal{H}_{\mathrm{cav}} \otimes \mathcal{H}_{\mathrm{atom}}$
with dimension $2N_{\mathrm{Fock}}$.

\textbf{Basis convention.} A basis state $\ket{n, s}$ is indexed as
$i = 2n + s$, where $n \in \{0,\ldots,N_F - 1\}$ is the photon number
and $s \in \{0 = e, 1 = g\}$ is the atomic spin.

\textbf{Hamiltonian (on resonance, $\omega_c = \omega_a = \omega$):}
\begin{equation}
  H = \omega\,a^\dagger a
    + \frac{\omega}{2}\,\sigma_z
    + g\,(a^\dagger\sigma_- + a\,\sigma_+)
  \label{eq:jc-ham-full}
\end{equation}

\textbf{Conservation law.} The operator $\hat{N} = a^\dagger a + \sigma_+\sigma_-$
is conserved by $H$. Starting from $\ket{1, e}$ ($\hat{N} = 2$),
the system oscillates within the two-state manifold
$\{\ket{1,e}, \ket{2,g}\}$ — not between $\ket{1,e}$ and $\ket{0,g}$.
This is an important subtlety verified by the test below.

\begin{minted}[linenos]{cpp}
#include "liquid/liquid.hpp"
#include <cmath>
#include <cstdio>
#include <vector>
using namespace liquid;
using namespace liquid::trajectory;
using namespace liquid::ode;

// Basis: |n, s> -> index 2*n + s  (s=0: excited, s=1: ground)
static Idx jc_idx(int n, int s) { return static_cast<Idx>(2 * n + s); }

// Cavity photon loss: L1 = sqrt(kappa) * a
// a|n,s> = sqrt(n)|n-1,s>
static SparseOperator make_L_cavity(int N_fock, double kappa) {
    const Dim dim = 2 * N_fock;
    std::vector<Triplet> trips;
    for (int n = 1; n < N_fock; ++n) {
        const double v = std::sqrt(kappa * static_cast<double>(n));
        trips.push_back({jc_idx(n-1, 0), jc_idx(n, 0), Scalar{v, 0.0}});
        trips.push_back({jc_idx(n-1, 1), jc_idx(n, 1), Scalar{v, 0.0}});
    }
    return SparseOperator(dim, std::move(trips));
}

// Atomic decay: L2 = sqrt(gamma_a) * sigma_minus
// sigma_minus|n,e> = |n,g>
static SparseOperator make_L_atomic(int N_fock, double gamma_a) {
    const Dim dim = 2 * N_fock;
    std::vector<Triplet> trips;
    const double v = std::sqrt(gamma_a);
    for (int n = 0; n < N_fock; ++n)
        trips.push_back({jc_idx(n, 1), jc_idx(n, 0), Scalar{v, 0.0}});
    return SparseOperator(dim, std::move(trips));
}

// JC Hamiltonian: omega*adaga + (omega/2)*sigma_z + g*(adagsigma_- + a*sigma_+)
// Coupling term: adagsigma_-|n,e> = sqrt(n+1)|n+1,g>
static SparseOperator make_H_jc(int N_fock, double omega, double g) {
    const Dim dim = 2 * N_fock;
    std::vector<Triplet> trips;

    // Free cavity: omega * adaga diagonal
    for (int n = 0; n < N_fock; ++n) {
        const double diag = omega * static_cast<double>(n);
        trips.push_back({jc_idx(n,0), jc_idx(n,0), Scalar{diag, 0.0}});
        trips.push_back({jc_idx(n,1), jc_idx(n,1), Scalar{diag, 0.0}});
    }

    // Free atom: (omega/2)*sigma_z diagonal
    for (int n = 0; n < N_fock; ++n) {
        trips.push_back({jc_idx(n,0), jc_idx(n,0), Scalar{ omega/2.0, 0.0}});
        trips.push_back({jc_idx(n,1), jc_idx(n,1), Scalar{-omega/2.0, 0.0}});
    }

    // JC interaction: adagsigma_- couples |n,e> -> |n+1,g>
    for (int n = 0; n < N_fock - 1; ++n) {
        const double v = g * std::sqrt(static_cast<double>(n + 1));
        trips.push_back({jc_idx(n+1, 1), jc_idx(n, 0), Scalar{v, 0.0}});  // adagsigma_-
        trips.push_back({jc_idx(n, 0), jc_idx(n+1, 1), Scalar{v, 0.0}});  // h.c. a*sigma_+
    }
    return SparseOperator(dim, std::move(trips));
}

// Photon number observable: n_hat = sum_n n * (|n,e><n,e| + |n,g><n,g|)
static ObservableDef make_photon_number_obs(int N_fock) {
    const Dim dim = 2 * N_fock;
    std::vector<Triplet> trips;
    for (int n = 0; n < N_fock; ++n) {
        const double val = static_cast<double>(n);
        trips.push_back({jc_idx(n,0), jc_idx(n,0), Scalar{val, 0.0}});
        trips.push_back({jc_idx(n,1), jc_idx(n,1), Scalar{val, 0.0}});
    }
    SparseOperator n_op(dim, std::move(trips));
    return {"photon_number",
            [op = std::move(n_op)](const StateVector& psi) -> Real {
                const Real ns = psi.norm_sq();
                return ns > 1e-30
                    ? sparse_expectation(op, psi).real() / ns
                    : 0.0;
            }};
}

// Excited-state population observable
static ObservableDef make_excited_pop_obs(int N_fock) {
    const Dim dim = 2 * N_fock;
    std::vector<Triplet> trips;
    for (int n = 0; n < N_fock; ++n)
        trips.push_back({jc_idx(n,0), jc_idx(n,0), Scalar{1.0, 0.0}});
    SparseOperator pe_op(dim, std::move(trips));
    return {"excited_pop",
            [op = std::move(pe_op)](const StateVector& psi) -> Real {
                const Real ns = psi.norm_sq();
                return ns > 1e-30
                    ? sparse_expectation(op, psi).real() / ns
                    : 0.0;
            }};
}

int main() {
    const int    N_fock   = 12;    // Fock space truncation
    const double omega    = 1.0;   // cavity / atom frequency
    const double g        = 0.5;   // vacuum Rabi coupling
    const double kappa    = 0.2;   // cavity photon loss rate
    const double gamma_a  = 0.05;  // atomic decay rate
    const double T        = 5.0;   // simulation time

    // Build sparse JC system
    SparseOperator H   = make_H_jc(N_fock, omega, g);
    SparseOperator L1  = make_L_cavity(N_fock, kappa);
    SparseOperator L2  = make_L_atomic(N_fock, gamma_a);

    std::vector<SparseOperator> ops;
    ops.push_back(std::move(L1));
    ops.push_back(std::move(L2));

    SparseOpenSystem sys(
        std::move(H),
        LindbladSet<SparseTag>(std::move(ops)));

    // Observables
    std::vector<ObservableDef> obs_defs;
    obs_defs.push_back(make_photon_number_obs(N_fock));
    obs_defs.push_back(make_excited_pop_obs(N_fock));

    // Stopping criteria
    StoppingCriteria sc;
    sc.min_trajectories = 500;
    sc.max_trajectories = 5000;
    sc.target_rel_sem   = 0.03;

    EnsembleConfig ec;
    ec.global_seed           = 42;
    ec.diag_level            = DiagnosticLevel::None;
    ec.propagator.dt_initial = 5e-3;

    // Use the sparse+DOPRI45 factory (SparseOpenSystem requires
    // Simulation::make_sparse_dopri45; SimulationBuilder supports
    // DenseOperator only in v0.7.0)
    Simulation sim = Simulation::make_sparse_dopri45(
        std::move(sys), std::move(obs_defs), sc, ec);

    // Initial state: |n=1, e> -- one photon, atom excited
    StateVector psi0(2 * N_fock);
    psi0[jc_idx(1, 0)] = Scalar{1.0, 0.0};

    auto result = sim.run(psi0, 0.0, T);

    std::printf("=== Jaynes-Cummings (N_fock=%d, g/kappa=%.1f) ===\n",
        N_fock, g / kappa);
    std::printf("  <n_photon>  = %.4f +/- %.4f  (N=%zu)\n",
        result.observables[0].mean, result.observables[0].sem,
        result.total_trajectories);
    std::printf("  <rho_e>     = %.4f +/- %.4f\n",
        result.observables[1].mean, result.observables[1].sem);
    std::printf("  Mean jumps  = %.2f per trajectory\n",
        result.mean_jumps_per_trajectory);
    std::printf("  Convergence = %s\n",
        result.convergence.reason.c_str());

    return 0;
}
\end{minted}

\begin{notebox}
  The driven-dissipative Jaynes--Cummings parameter sweep (sweeping $g/\kappa$
  to observe photon blockade) is provided in full in
  \filename{examples/jc\_steady\_state.cpp}, using
  \code{ParameterSweepBuilder::parameter\_logrange()}.
\end{notebox}

% -----------------------------------------------------------------------------
\section{Example 4: Parameter Sweeps}
\label{sec:ex-sweep}
% -----------------------------------------------------------------------------

\code{ParameterSweepBuilder} runs the same simulation at a list of
parameter values and collects fully converged \code{EnsembleResult}
objects at each point. Two sweep modes are supported: linear
(\code{.parameter("name", \{v1, v2, ...\})}) and log-spaced
(\code{.parameter\_logrange("name", lo, hi, N)}).

The example below sweeps $\Omega/\gamma \in \{0.1, 0.5, 1.0, 2.0\}$
for the driven two-level system and verifies the steady-state formula
\cref{eq:driven-ss} at each point.

\begin{minted}[linenos]{cpp}
#include "liquid/liquid.hpp"
#include <cmath>
#include <cstdio>
using namespace liquid;
using namespace liquid::ensemble;

int main() {
    const double gamma = 1.0;
    const double T_ss  = 15.0;  // sufficient time to reach steady state

    // Initial state: excited |e>
    StateVector psi0(2);
    psi0[0] = Scalar{1.0, 0.0};

    // Analytic steady-state formula
    auto exact_ss = [&](double Omega) {
        const double r = Omega / gamma;
        return 4.0 * r * r / (1.0 + 8.0 * r * r);
    };

    // Define the sweep over Rabi frequency Omega
    auto sweep = ParameterSweepBuilder{}
        .parameter("Omega", {0.1, 0.5, 1.0, 2.0})
        .simulation_factory([gamma](double Omega) {
            // H = Omega * sigma_x  (resonant drive)
            DenseOperator H(2);
            H(0, 1) = Scalar{Omega, 0.0};
            H(1, 0) = Scalar{Omega, 0.0};

            // L = sqrt(gamma) * sigma_minus
            DenseOperator L(2);
            L(1, 0) = Scalar{std::sqrt(gamma), 0.0};

            // Observable: rho_ee = |e><e|
            DenseOperator rho_ee_op(2);
            rho_ee_op(0, 0) = Scalar{1.0, 0.0};

            return SimulationBuilder{}
                .hamiltonian(std::move(H))
                .collapse_operator(std::move(L))
                .observe("rho_ee", std::move(rho_ee_op))
                .seed(42)
                .dt(5e-4)
                .stop_at_sem(0.02)          // 2% relative SEM at each point
                .min_trajectories(500)
                .max_trajectories(5000)
                .build();
        })
        .initial_state(psi0)
        .time_interval(0.0, T_ss)
        .build();

    SweepResult result = sweep.run();

    // Print table with comparison to analytic formula
    std::printf("%-10s  %-12s  %-12s  %-10s  %-8s  %s\n",
        "Omega", "rho_ee(MCWF)", "rho_ee(exact)", "rel_err(%)", "N", "converged");
    std::printf("%s\n", std::string(72, '-').c_str());

    for (const auto& pt : result.points) {
        const auto& obs   = pt.result.observables[0];
        const double exact = exact_ss(pt.param_value);
        const double rel_err = 100.0 * std::abs(obs.mean - exact) / exact;

        std::printf("%-10.4f  %-12.4f  %-12.4f  %-10.2f  %-8zu  %s\n",
            pt.param_value,
            obs.mean,
            exact,
            rel_err,
            pt.result.total_trajectories,
            pt.result.convergence.reason.c_str());
    }

    // Save full results (CSV columns: param,obs,mean,sem,rel_sem,N,...)
    result.save_csv("sweep_Omega.csv");
    std::printf("\nSaved: sweep_Omega.csv\n");

    return 0;
}
\end{minted}

\begin{perfbox}
  When running sweeps with \code{num\_threads > 1}, the threads are
  used \emph{within} each simulation point (ensemble parallelism),
  not across sweep points. Sweep points are executed sequentially;
  this allows early stopping to adapt $N$ independently at each point
  rather than committing to a fixed $N$ before the sweep starts.
\end{perfbox}

% -----------------------------------------------------------------------------
\section{Example 5: Convergence-Based Stopping}
\label{sec:ex-stopping}
% -----------------------------------------------------------------------------

\code{StoppingCriteria} provides three independent stopping conditions:
SEM target, trajectory budget, and wall-clock budget. The
\code{ConvergenceMonitor} checks them after every batch of completed
trajectories. The example below demonstrates all three and verifies
that the SEM scales as $1/\sqrt{N}$.

\begin{minted}[linenos]{cpp}
#include "liquid/liquid.hpp"
#include <cmath>
#include <cstdio>
using namespace liquid;
using namespace liquid::ensemble;

static DenseOperator make_H(double omega) {
    DenseOperator H(2);
    H(0,0) = Scalar{ omega/2.0, 0.0};
    H(1,1) = Scalar{-omega/2.0, 0.0};
    return H;
}
static DenseOperator make_L(double gamma) {
    DenseOperator L(2);
    L(1,0) = Scalar{std::sqrt(gamma), 0.0};
    return L;
}
static DenseOperator make_sz() {
    DenseOperator sz(2);
    sz(0,0) = Scalar{1.0, 0.0};
    sz(1,1) = Scalar{-1.0, 0.0};
    return sz;
}

int main() {
    const double omega = 1.0;
    const double gamma = 1.0;
    const double T     = 1.0;

    StateVector psi0(2);
    psi0[0] = Scalar{1.0, 0.0};

    const double exact = 2.0 * std::exp(-gamma * T) - 1.0;

    // -- 1. SEM-based stopping (1% relative SEM target) ---------------------
    {
        Simulation sim = SimulationBuilder{}
            .hamiltonian(make_H(omega))
            .collapse_operator(make_L(gamma))
            .observe("sz", make_sz())
            .seed(42)
            .dt(1e-3)
            .stop_at_sem(0.01)
            .min_trajectories(50)
            .max_trajectories(50000)
            .build();

        auto result = sim.run(psi0, 0.0, T);
        const auto& obs = result.observables[0];

        std::printf("=== SEM-based stopping ===\n");
        std::printf("  <sz> = %.4f +/- %.4f  rel_sem = %.4f  N = %zu\n",
            obs.mean, obs.sem, obs.rel_sem, result.total_trajectories);
        std::printf("  Exact = %.4f  |error| = %.4f sigma\n",
            exact, std::abs(obs.mean - exact) / obs.sem);
        std::printf("  Decision: %s\n\n",
            result.convergence.reason.c_str());
    }

    // -- 2. Wall-clock budget (0.3 s limit, impossibly tight SEM) -----------
    {
        Simulation sim = SimulationBuilder{}
            .hamiltonian(make_H(omega))
            .collapse_operator(make_L(gamma))
            .observe("sz", make_sz())
            .seed(999)
            .dt(1e-3)
            .stop_at_sem(1e-9)          // impossible: triggers budget
            .min_trajectories(10)
            .max_trajectories(1000000)
            .max_wall_time(0.3)         // 300 ms hard limit
            .build();

        auto result = sim.run(psi0, 0.0, T);

        std::printf("=== Wall-clock budget ===\n");
        std::printf("  Decision  : %s\n",
            result.convergence.reason.c_str());
        std::printf("  N achieved: %zu  wall = %.3f s\n\n",
            result.total_trajectories,
            result.total_wall_time_seconds);
    }

    // -- 3. SEM scales as 1/sqrt(N): verify with three ensemble sizes --------
    {
        const std::size_t sizes[3] = {100, 400, 1600};
        double sems[3];

        std::printf("=== SEM scaling: expected SEM ~ 1/sqrt(N) ===\n");
        for (int k = 0; k < 3; ++k) {
            Simulation sim = SimulationBuilder{}
                .hamiltonian(make_H(omega))
                .collapse_operator(make_L(gamma))
                .observe("sz", make_sz())
                .seed(static_cast<Seed>(k * 777 + 42))
                .dt(1e-3)
                .min_trajectories(sizes[k])
                .max_trajectories(sizes[k])
                .build();

            auto result  = sim.run(psi0, 0.0, T);
            sems[k]      = result.observables[0].sem;
            std::printf("  N = %4zu  SEM = %.5f\n", sizes[k], sems[k]);
        }

        // Ratios should be ~2 (since N quadruples each time)
        std::printf("  SEM ratio N=100/400  = %.3f  (expected ~2.0)\n",
            sems[0] / sems[1]);
        std::printf("  SEM ratio N=400/1600 = %.3f  (expected ~2.0)\n",
            sems[1] / sems[2]);
    }

    return 0;
}
\end{minted}

\begin{warnbox}
  \code{max\_wall\_time} measures elapsed real time from the first
  \code{sim.run()} call. It does \emph{not} account for operator
  construction or build time. On heavily loaded machines, the wall
  clock continues to tick even while the process is descheduled.
  Use \code{max\_trajectories} as the primary safety ceiling;
  \code{max\_wall\_time} is a soft backstop for interactive sessions.
\end{warnbox}

% -----------------------------------------------------------------------------
\section{Example 6: Sparse Hilbert Spaces}
\label{sec:ex-sparse}
% -----------------------------------------------------------------------------

This example illustrates the three design patterns for sparse
Hilbert-space simulations: (i)~using \code{SparseOperator} directly
for a high-dimensional bosonic mode, (ii)~exploiting the
\code{Simulation::make\_sparse\_dopri45} factory for non-builder-compatible
system types, and (iii)~verifying the sparse--dense equivalence
for a small test system.

\subsubsection*{Pattern A: Bosonic Fock space (pure sparse)}

For a single driven, damped harmonic oscillator (cavity without atom):

\begin{equation}
  H = \Delta\,a^\dagger a + \varepsilon(a + a^\dagger),
  \qquad
  L = \sqrt{\kappa}\,a
  \label{eq:harmonic-osc}
\end{equation}

The steady-state mean photon number has the exact value
$\langle a^\dagger a\rangle_{\mathrm{ss}} = \varepsilon^2/\kappa^2$
in the linear regime ($\varepsilon \ll \kappa$).

\begin{minted}[linenos]{cpp}
#include "liquid/liquid.hpp"
#include <cmath>
#include <cstdio>
#include <vector>
using namespace liquid;
using namespace liquid::ensemble;

// Build the annihilation operator a for an N_fock-dimensional Fock space.
// a|n> = sqrt(n)|n-1>  for n >= 1; a|0> = 0
static SparseOperator make_a_fock(Dim N_fock) {
    std::vector<Triplet> trips;
    trips.reserve(N_fock - 1);
    for (Dim n = 1; n < N_fock; ++n) {
        trips.push_back({
            n - 1,                            // row (output state index)
            n,                                // col (input state index)
            Scalar{std::sqrt(static_cast<double>(n)), 0.0}});
    }
    return SparseOperator(N_fock, std::move(trips));
}

// Driven harmonic oscillator Hamiltonian:
// H = Delta * adaga + epsilon * (a + adag)
static SparseOperator make_H_harmonic(Dim N_fock, double Delta, double eps) {
    std::vector<Triplet> trips;
    // Diagonal: Delta * n
    for (Dim n = 0; n < N_fock; ++n)
        trips.push_back({n, n, Scalar{Delta * static_cast<double>(n), 0.0}});
    // Drive: epsilon * a  (lower triangular)
    for (Dim n = 1; n < N_fock; ++n)
        trips.push_back({n-1, n, Scalar{eps * std::sqrt(static_cast<double>(n)), 0.0}});
    // Drive: epsilon * adag  (upper triangular)
    for (Dim n = 1; n < N_fock; ++n)
        trips.push_back({n, n-1, Scalar{eps * std::sqrt(static_cast<double>(n)), 0.0}});
    return SparseOperator(N_fock, std::move(trips));
}

// Cavity decay: L = sqrt(kappa) * a
static SparseOperator make_L_harmonic(Dim N_fock, double kappa) {
    std::vector<Triplet> trips;
    for (Dim n = 1; n < N_fock; ++n)
        trips.push_back({
            n - 1, n,
            Scalar{std::sqrt(kappa * static_cast<double>(n)), 0.0}});
    return SparseOperator(N_fock, std::move(trips));
}

// Photon number observable: adaga = diag(0, 1, 2, ..., N_fock-1)
static ObservableDef make_n_obs(Dim N_fock) {
    std::vector<Triplet> trips;
    for (Dim n = 0; n < N_fock; ++n)
        trips.push_back({n, n, Scalar{static_cast<double>(n), 0.0}});
    SparseOperator n_op(N_fock, std::move(trips));
    return {"photon_n",
            [op = std::move(n_op)](const StateVector& psi) -> Real {
                const Real ns = psi.norm_sq();
                return ns > 1e-30
                    ? sparse_expectation(op, psi).real() / ns
                    : 0.0;
            }};
}

int main() {
    const Dim    N_fock = 40;    // truncation: must satisfy N >> <n>_ss
    const double Delta  = 0.0;   // on resonance with drive
    const double kappa  = 1.0;   // cavity decay rate
    const double eps    = 0.2;   // drive amplitude (weak: eps << kappa)
    const double T_ss   = 15.0;  // run until steady state

    // Exact steady-state photon number (linear cavity)
    const double exact_n_ss = (eps * eps) / (kappa * kappa);

    std::printf("N_fock=%zu  Delta=%.1f  kappa=%.1f  eps=%.2f\n",
        N_fock, Delta, kappa, eps);
    std::printf("Exact <n>_ss = eps^2/kappa^2 = %.4f\n\n", exact_n_ss);

    // Verify truncation is adequate: <n>_ss << N_fock
    if (exact_n_ss > static_cast<double>(N_fock) / 2.0) {
        std::fprintf(stderr,
            "WARNING: N_fock=%zu may be too small for <n>_ss=%.1f\n",
            N_fock, exact_n_ss);
    }

    // Build sparse system
    SparseOperator H  = make_H_harmonic(N_fock, Delta, eps);
    SparseOperator L  = make_L_harmonic(N_fock, kappa);

    std::vector<SparseOperator> ops;
    ops.push_back(std::move(L));
    SparseOpenSystem sys(std::move(H), LindbladSet<SparseTag>(std::move(ops)));

    // Report sparsity advantage over dense
    const double nnz      = static_cast<double>(sys.H_eff().nnz());
    const double dense_sq = static_cast<double>(N_fock * N_fock);
    std::printf("H_eff sparsity: %.2f%%  (%g non-zeros vs %g dense)\n\n",
        100.0 * nnz / dense_sq, nnz, dense_sq);

    std::vector<ObservableDef> obs_defs;
    obs_defs.push_back(make_n_obs(N_fock));

    StoppingCriteria sc;
    sc.min_trajectories = 300;
    sc.max_trajectories = 5000;
    sc.target_rel_sem   = 0.02;

    EnsembleConfig ec;
    ec.global_seed           = 42;
    ec.diag_level            = DiagnosticLevel::None;
    ec.propagator.dt_initial = 5e-3;

    // Factory for SparseOpenSystem + DOPRI45
    Simulation sim = Simulation::make_sparse_dopri45(
        std::move(sys), std::move(obs_defs), sc, ec);

    // Initial state: vacuum |0>
    StateVector psi0(N_fock);
    psi0[0] = Scalar{1.0, 0.0};

    auto result = sim.run(psi0, 0.0, T_ss);

    const auto& obs = result.observables[0];
    std::printf("=== Driven harmonic oscillator (N_fock=%zu) ===\n", N_fock);
    std::printf("  <n>(MCWF)  = %.4f +/- %.4f  (N=%zu)\n",
        obs.mean, obs.sem, result.total_trajectories);
    std::printf("  <n>(exact) = %.4f\n", exact_n_ss);
    std::printf("  |error|    = %.4f  (%.1f sigma)\n",
        std::abs(obs.mean - exact_n_ss),
        std::abs(obs.mean - exact_n_ss) / obs.sem);

    return 0;
}
\end{minted}

\begin{perfbox}
  The Hamiltonian \cref{eq:harmonic-osc} has $3N_F - 2$ non-zero elements
  (one diagonal plus two off-diagonals of length $N_F - 1$ each).
  For $N_F = 40$, this is 118 non-zeros versus $1600$ for a dense
  matrix --- a $13.6\times$ memory saving and a comparable
  reduction in matrix-vector product cost. The savings grow linearly
  with $N_F$: for $N_F = 200$, the ratio exceeds $66\times$.
\end{perfbox}

% ============================================================
%  PART IV — PERFORMANCE AND EXTENSION
% ============================================================
\part{Performance and Extension}

% -------------------------------------------------------------
\chapter{Performance Guide}
\label{chap:performance}
% -------------------------------------------------------------

\section{Computational Complexity Model}
\label{sec:complexity}

\begin{table}[H]
\centering
\caption{Computational cost of each layer per trajectory.
         $N$ = Hilbert dimension, $K$ = jump operators, $M$ = ODE steps,
         $J$ = number of jumps.}
\label{tab:complexity}
\begin{tabularx}{\linewidth}{L{30mm}L{25mm}X}
\toprule
\textbf{Operation} & \textbf{Cost} & \textbf{Notes} \\
\midrule
Dense \code{apply\_add}  & $\Order(N^2)$  & Dominates for $N > 32$ \\
Sparse \code{apply\_add} & $\Order(\mathrm{nnz})$ & $\ll N^2$ for sparse systems \\
RK4 step                 & $4 \cdot \Order(N^2)$ & 4 RHS evaluations \\
DOPRI45 step (accepted)  & $6 \cdot \Order(N^2)$ & 6 RHS (FSAL saves one) \\
Jump detection           & $\Order(1)$    & Norm comparison only \\
Jump channel selection   & $\Order(KN^2)$ & $K$ matvecs to compute $p_k$ \\
State renormalization    & $\Order(N)$    & Linear sweep \\
Observable evaluation    & $\Order(N^2)$  & Once per trajectory \\
\midrule
\textbf{Per-trajectory total} & $\Order(M \cdot N^2)$ & Dominated by ODE steps \\
\textbf{Ensemble total}       & $\Order(M \cdot N^2 \cdot T_{\mathrm{traj}})$ & Linear in $T$ \\
\bottomrule
\end{tabularx}
\end{table}

\section{Memory Hierarchy Considerations}
\label{sec:memory-hierarchy}

\begin{perfbox}
  For $N \leq 64$, the wavefunction $\ket{\psi}$ fits in
  $64 \times 16 = 1\,\mathrm{kB}$, which resides in L1 cache.
  The RK4 scratch buffers (5 vectors $= 5\,\mathrm{kB}$) fit in L1/L2.
  A $64 \times 64$ dense matrix ($= 64\,\mathrm{kB}$) fits in L2.
  Keeping all working data in L1/L2 is the primary performance objective
  for small systems. Use \code{-march=native} to enable AVX2
  auto-vectorization of the \code{add\_scaled} and \code{apply\_add} loops.
\end{perfbox}

\section{Dense vs.\ Sparse Crossover}
\label{sec:dense-sparse-crossover}

\begin{figure}[H]
\centering
\begin{tikzpicture}
\begin{axis}[
  width=0.75\linewidth, height=50mm,
  xlabel={Hilbert space dimension $N$},
  ylabel={Relative speedup (sparse/dense)},
  xmin=4, xmax=512,
  ymin=0, ymax=20,
  xmode=log, log basis x=2,
  grid=major,
  grid style={dotted, gray!50},
  xtick={4,8,16,32,64,128,256,512},
  xticklabels={4,8,16,32,64,128,256,512},
  legend pos=north west,
  title={Sparse speedup for bosonic annihilation operator $a$},
]
% nnz = N-1 for bosonic ladder; dense = N^2
% speedup = N^2 / (N-1) ~= N
\addplot[thick, color=LiquidTeal, mark=*]
  coordinates {(4,1.3)(8,2.3)(16,4.3)(32,8.3)(64,16.3)(128,32.3)};
\addlegendentry{$\approx N$ speedup ($\mathrm{nnz} = N-1$)}
\addplot[thick, dashed, color=LiquidGold]
  coordinates {(4,1)(512,1)};
\addlegendentry{Dense baseline}
\end{axis}
\end{tikzpicture}
\caption{Theoretical speedup of sparse vs.\ dense matrix-vector product
         for the bosonic annihilation operator $a$ (non-zero elements: $N-1$).
         The crossover to practical benefit occurs around $N \approx 32$.}
\label{fig:sparse-speedup}
\end{figure}

\section{Multi-Threading Scaling}
\label{sec:threading}

The ensemble is \emph{embarrassingly parallel}: each trajectory is
independent. With OpenMP enabled:

\begin{minted}{cpp}
// In EnsembleConfig:
config.num_threads = std::thread::hardware_concurrency();
\end{minted}

Each thread owns:
\begin{itemize}
  \item One \code{MCWFPropagator} (and its scratch buffers).
  \item One \code{ODE stepper} (and its scratch buffers).
  \item One \code{ObservableAccumulator} (merged after each batch).
\end{itemize}

The shared read-only objects — \code{OpenSystem}, all operator matrices,
observable definitions — are accessed concurrently without locks because
they are immutable after construction.

\begin{warnbox}
  OpenMP must be enabled at compile time (\code{-fopenmp}) and at
  CMake configuration time (\code{-DLIQUID\_OPENMP=ON}).
  Setting \code{config.num\_threads > 1} on a build without OpenMP
  has no effect; the simulation runs serially.
\end{warnbox}

% -------------------------------------------------------------
\chapter{Extending LiQuID}
\label{chap:extending}
% -------------------------------------------------------------

\section{Custom ODE Stepper}
\label{sec:custom-stepper}

Any class satisfying the Stepper concept (\cref{sec:stepper-concept})
can be used as a drop-in replacement:

\begin{minted}[linenos]{cpp}
// Minimal custom stepper: symplectic Euler (first-order)
class SymplecticEulerStepper {
public:
    explicit SymplecticEulerStepper(Real dt) : dt_(dt) {}

    template<typename RHSFn>
    ode::StepResult try_step(Real t,
                              const StateVector& psi_in,
                              StateVector& psi_out,
                              Real dt_suggested,
                              Real t_max,
                              RHSFn&& rhs_fn)
    {
        const Real dt = std::min({dt_suggested, dt_, t_max - t});
        ensure_scratch(psi_in.size());

        // Compute derivative: k = f(t, psi_in)
        k_.set_zero();
        rhs_fn(t, psi_in, k_);

        // Update: psi_out = psi_in + dt * k
        psi_out.copy_from(psi_in);
        psi_out.add_scaled(k_, Scalar{dt, 0.0});

        return ode::StepResult{true, dt, dt_, 0.0};
    }

    void reset(Real dt_initial) noexcept { dt_ = dt_initial; }

private:
    void ensure_scratch(Dim n) {
        if (k_.size() != n) k_ = StateVector(n);
    }

    Real        dt_;
    StateVector k_;
};

// Use it with MCWFPropagator:
MCWFPropagator<DenseOpenSystem, SymplecticEulerStepper> prop(
    &sys,
    SymplecticEulerStepper(1e-4),
    cfg);
\end{minted}

\section{Custom Allocation Policy}
\label{sec:custom-policy}

\begin{minted}[linenos]{cpp}
// Policy that preferentially runs more trajectories when
// the failure rate is below 5%, and falls back to uniform otherwise.
AllocatorPolicy my_policy = [seed = 42ULL](
    const EnsembleSummary& summary) mutable -> AllocationDecision
{
    const double failure_rate =
        static_cast<double>(summary.trajectories_failed) /
        std::max(summary.trajectories_completed, std::size_t{1});

    AllocationDecision dec;
    dec.action = AllocationAction::SpawnNew;

    if (failure_rate < 0.05) {
        // Low failure rate: use a deterministic seed sequence
        dec.seed = seed++;
    } else {
        // High failure rate: fall back to random seed
        dec.seed = std::hash<double>{}(summary.wall_time_elapsed);
    }
    return dec;
};

// Pass to Simulation::make() or EnsembleManager
auto sim = Simulation::make<DenseOpenSystem, ode::RK4Stepper>(
    std::move(system),
    std::move(observable_defs),
    stopping_criteria,
    ensemble_config,
    my_policy);
\end{minted}

% -------------------------------------------------------------
\chapter{Test Suite}
\label{chap:tests}
% -------------------------------------------------------------

\section{Unit Tests}
\label{sec:unit-tests}

\begin{table}[H]
\centering
\caption{Unit test executables and their scope.}
\label{tab:unit-tests}
\begin{tabularx}{\linewidth}{L{48mm}X}
\toprule
\textbf{Target} & \textbf{Scope} \\
\midrule
\code{test\_rng}              & xoshiro256** uniformity, period, seed independence \\
\code{test\_linalg}           & \code{StateVector}, \code{DenseOperator}: all methods and edge cases \\
\code{test\_sparse}           & CSR construction, \code{apply\_add}, duplicate triplet handling \\
\code{test\_opensystem}       & \code{H\_eff} precomputation, jump probabilities, apply\_jump \\
\code{test\_trajectory\_state}& Factory postconditions, all seven invariants \\
\code{test\_running\_statistics}& Welford online mean/variance, single-sample edge case \\
\code{test\_ensemble}         & Convergence monitor, observable accumulator, parallel merge \\
\code{test\_dopri45}          & FSAL correctness, PI controller, step rejection \\
\code{test\_ml\_policy}       & Feature extraction, MLP forward pass, SGD convergence \\
\code{test\_parameter\_sweep} & CSV output, factory invocation, OpenMP sweep \\
\bottomrule
\end{tabularx}
\end{table}

\section{Validation Tests}
\label{sec:validation-tests}

\begin{table}[H]
\centering
\caption{Physics validation tests and their analytical references.}
\label{tab:validation-tests}
\begin{tabularx}{\linewidth}{L{42mm}X}
\toprule
\textbf{Target} & \textbf{Reference} \\
\midrule
\code{val\_two\_level\_decay}
  & $\langle\sigma_z\rangle = 2e^{-\gamma t}-1$; passes within $3\sigma$ \\
\code{val\_driven\_two\_level}
  & Driven qubit: steady-state Bloch vector vs.\ optical Bloch equations \\
\code{val\_ensemble\_convergence}
  & SEM $\propto 1/\sqrt{N}$: fit slope $= -0.5 \pm 0.05$ \\
\code{val\_cavity\_qed}
  & Jaynes--Cummings photon statistics: $g^{(2)}(0) < 1$ (antibunching) \\
\code{val\_smart\_stopping}
  & Early stopping triggers before \code{max\_trajectories} when converged \\
\code{val\_ml\_policy}
  & ML policy SEM reduction rate $\geq$ uniform policy within $10\%$ \\
\code{val\_parameter\_sweep}
  & Sweep monotonicity: $\langle\sigma_z\rangle_{\mathrm{ss}}$ decreasing with $\gamma$ \\
\bottomrule
\end{tabularx}
\end{table}

Run all tests with:
\begin{minted}{bash}
ctest --test-dir build --output-on-failure -j$(nproc)
\end{minted}

% -------------------------------------------------------------
\chapter{Known Limitations and Roadmap}
\label{chap:limitations}
% -------------------------------------------------------------

\section{Known Limitations in v0.7.0}
\label{sec:known-limitations}

\begin{table}[H]
\centering
\caption{Known limitations and planned resolution phases.}
\label{tab:limitations}
\begin{tabularx}{\linewidth}{L{55mm}L{15mm}X}
\toprule
\textbf{Limitation} & \textbf{Phase} & \textbf{Notes} \\
\midrule
Phase 1 jump location: no bisection for RK4 & Phase 3 done & DOPRI45 path uses bisection \\
No BLAS dependency & Phase 3+ & Manual loops; switch to BLAS for $N > 100$ \\
\code{SimulationBuilder}: dense operators only & Phase 4 & Use \code{Simulation::make()} for sparse \\
No Python bindings & Post-1.0 & pybind11 wrapper planned \\
No GPU backend & Post-1.0 & Explicit deferred decision \\
ML policy: seeds only, no rare-event biasing & Phase 6 & Infrastructure present, biasing deferred \\
\bottomrule
\end{tabularx}
\end{table}

% -------------------------------------------------------------
%  APPENDICES
% -------------------------------------------------------------
\appendix

\chapter{Mathematical Reference}
\label{app:math}

\section{Lindblad vs.\ GKSL Notation}
\label{app:notation}

Different communities use different conventions for \cref{eq:lindblad}.
LiQuID uses the following (natural units $\hbar = 1$):

\begin{align}
  \frac{d\densmat}{dt}
    &= -i[H, \densmat]
     + \sum_k \left(L_k\densmat L_k^\dagger
       - \tfrac{1}{2}L_k^\dagger L_k\densmat
       - \tfrac{1}{2}\densmat L_k^\dagger L_k\right) \notag\\
    &\equiv -i[H, \densmat]
     + \sum_k \mathcal{D}[L_k]\densmat
  \label{eq:lindblad-dissipator}
\end{align}

where $\mathcal{D}[L]\densmat \equiv L\densmat L^\dagger - \tfrac{1}{2}\{L^\dagger L, \densmat\}$.
Some references write $\gamma_k L_k$ with unit-norm $L_k$; others absorb
$\sqrt{\gamma_k}$ into $L_k$. LiQuID uses the latter convention:
$L_k = \sqrt{\gamma_k}\,\tilde{L}_k$ where $\tilde{L}_k$ is dimensionless.

\section{Standard Error of the Mean}
\label{app:sem}

For $M$ independent MCWF samples $x_1, \ldots, x_M$ of an observable,
LiQuID computes:

\begin{align}
  \hat{\mu}   &= \frac{1}{M}\sum_{m=1}^M x_m \\
  \hat{s}^2   &= \frac{1}{M-1}\sum_{m=1}^M (x_m - \hat{\mu})^2 \\
  \widehat{\mathrm{SEM}} &= \frac{\hat{s}}{\sqrt{M}} \\
  \varepsilon_{\mathrm{rel}} &= \frac{\widehat{\mathrm{SEM}}}{|\hat{\mu}|}
\end{align}

The stopping criterion $\varepsilon_{\mathrm{rel}} < \varepsilon_{\mathrm{target}}$
is applied to the worst-case observable across all channels.
The variance $\hat{s}^2$ is computed via Welford's online algorithm
(implemented in \code{RunningStatistics}) to avoid catastrophic cancellation.

\chapter{API Quick Reference}
\label{app:api-ref}

\section{Builder Method Summary}
\label{app:builder-summary}

\begin{longtable}{L{42mm}L{28mm}L{38mm}}
\caption{Complete \code{SimulationBuilder} method reference.}
\label{tab:builder-ref}\\
\toprule
\textbf{Method} & \textbf{Type} & \textbf{Sets field} \\
\midrule
\endfirsthead
\midrule
\textbf{Method} & \textbf{Type} & \textbf{Sets field} \\
\midrule
\endhead
\bottomrule
\endfoot
\code{.hamiltonian(H)}       & \code{DenseOperator} & \code{H\_} \\
\code{.collapse\_operator(L)} & \code{DenseOperator} & adds to \code{L\_list\_} \\
\code{.observe(name, O)}     & \code{string, DenseOperator} & adds to \code{obs\_} \\
\code{.observe(name, fn)}    & \code{string, ObservableFn} & adds to \code{obs\_} \\
\code{.dt(h)}                & \code{Real} & \code{cfg.dt\_initial} \\
\code{.solver(type)}         & \code{SolverType} & \code{solver\_} \\
\code{.atol(e)}              & \code{Real} & \code{cfg.atol} \\
\code{.rtol(e)}              & \code{Real} & \code{cfg.rtol} \\
\code{.seed(s)}              & \code{Seed} & \code{ens\_cfg.global\_seed} \\
\code{.num\_threads(n)}      & \code{size\_t} & \code{ens\_cfg.num\_threads} \\
\code{.stop\_at\_sem(e)}     & \code{Real} & \code{sc.target\_rel\_sem} \\
\code{.min\_trajectories(n)} & \code{size\_t} & \code{sc.min\_trajectories} \\
\code{.max\_trajectories(n)} & \code{size\_t} & \code{sc.max\_trajectories} \\
\code{.max\_wall\_time(s)}   & \code{Real} & \code{sc.max\_wall\_seconds} \\
\code{.diagnostic\_level(l)} & \code{DiagnosticLevel} & \code{ens\_cfg.diag\_level} \\
\code{.build()}              & $\to$ \code{Simulation} & validates + constructs \\
\end{longtable}

% -------------------------------------------------------------
%  BIBLIOGRAPHY
% -------------------------------------------------------------
\backmatter

\begin{thebibliography}{99}
\bibitem{dalibard1992}
  J.~Dalibard, Y.~Castin, and K.~M\o{}lmer,
  ``Wave-function approach to dissipative processes in quantum optics,''
  \emph{Phys.\ Rev.\ Lett.}\ \textbf{68}, 580 (1992).

\bibitem{dum1992}
  R.~Dum, P.~Zoller, and H.~Ritsch,
  ``Monte Carlo simulation of the atomic master equation for spontaneous emission,''
  \emph{Phys.\ Rev.\ A}\ \textbf{45}, 4879 (1992).

\bibitem{molmer1993}
  K.~M\o{}lmer, Y.~Castin, and J.~Dalibard,
  ``Monte Carlo wave-function method in quantum optics,''
  \emph{J.\ Opt.\ Soc.\ Am.\ B}\ \textbf{10}, 524 (1993).

\bibitem{carmichael1993}
  H.~Carmichael,
  \emph{An Open Systems Approach to Quantum Optics}.
  Springer, Berlin, 1993.

\bibitem{gardiner2004}
  C.~W.~Gardiner and P.~Zoller,
  \emph{Quantum Noise}, 3rd ed.
  Springer, Berlin, 2004.

\bibitem{dormand1980}
  J.~R.~Dormand and P.~J.~Prince,
  ``A family of embedded Runge-Kutta formulae,''
  \emph{J.\ Comput.\ Appl.\ Math.}\ \textbf{6}, 19--26 (1980).

\bibitem{hairer1993}
  E.~Hairer, S.~P.~N\o{}rsett, and G.~Wanner,
  \emph{Solving Ordinary Differential Equations I}, 2nd ed.
  Springer, Berlin, 1993.

\bibitem{gustafsson1991}
  K.~Gustafsson,
  ``Control theoretic techniques for stepsize selection in explicit Runge-Kutta methods,''
  \emph{ACM Trans.\ Math.\ Softw.}\ \textbf{17}, 533--554 (1991).

\bibitem{blackman2018}
  D.~Blackman and S.~Vigna,
  ``Scrambled linear pseudorandom number generators,''
  \emph{ACM Trans.\ Math.\ Softw.}\ \textbf{47}, 1--32 (2021).
  \texttt{arXiv:1805.01407}.

\bibitem{welford1962}
  B.~P.~Welford,
  ``Note on a method for calculating corrected sums of squares and products,''
  \emph{Technometrics}\ \textbf{4}, 419--420 (1962).

\bibitem{lindblad1976}
  G.~Lindblad,
  ``On the generators of quantum dynamical semigroups,''
  \emph{Commun.\ Math.\ Phys.}\ \textbf{48}, 119--130 (1976).

\bibitem{gorini1976}
  V.~Gorini, A.~Kossakowski, and E.~C.~G.~Sudarshan,
  ``Completely positive dynamical semigroups of $N$-level systems,''
  \emph{J.\ Math.\ Phys.}\ \textbf{17}, 821--825 (1976).
\end{thebibliography}

% -------------------------------------------------------------
%  INDEX
% -------------------------------------------------------------
% Index omitted in this build (requires makeindex run)
% Run: makeindex liquid_manual.idx && pdflatex liquid_manual.tex

\end{document}
